\documentclass[12pt, a4paper]{article}
\usepackage{amsmath, amsthm, amssymb, amsfonts}
\usepackage{mathtools}
\usepackage[margin=2.5cm]{geometry}
\usepackage{enumitem}
\usepackage[normalem]{ulem}
\usepackage{hyperref}
% microtype removed (font expansion incompatibility)

\theoremstyle{definition}
\newtheorem{definition}{Definition}[section]
\newtheorem{example}{Example}[section]
\newtheorem{remark}{Remark}[section]

\theoremstyle{plain}
\newtheorem{theorem}{Theorem}[section]
\newtheorem{lemma}{Lemma}[section]
\newtheorem{proposition}{Proposition}[section]
\newtheorem{corollary}{Corollary}[section]

\DeclareMathOperator{\Exp}{Exp}
\DeclareMathOperator{\Poi}{Poi}
\DeclareMathOperator{\Bin}{Bin}
\newcommand{\PP}{\mathbb{P}}
\newcommand{\EE}{\mathbb{E}}
\newcommand{\FF}{\mathcal{F}}
\newcommand{\II}{\mathcal{I}}
\newcommand{\RR}{\mathbb{R}}
\newcommand{\NN}{\mathbb{N}}

\title{IE 622: Probability and Stochastic Processes\\
\large Lecture Notes (Transcribed)}
\date{IIT Bombay}

\begin{document}
\maketitle
\tableofcontents
\newpage

\section{Probability Foundations}

\begin{example}[Discrete probability space]
Tossing a coin, rolling a die.

\textbf{Discrete probability space:} $\Omega$ is a finite or countable sample space. For $A \subseteq \Omega$,
\begin{equation*}
  \PP(A) = \sum_{\omega \in A} p_\omega.
\end{equation*}

Toss a coin with bias $p$, $n$ times. Then $\Omega = \{0,1\}^n$, $\omega = (\omega_1, \omega_2, \ldots, \omega_n)$,
\begin{equation*}
  p_\omega = p^{\omega_1+\omega_2+\cdots+\omega_n}(1-p)^{n-(\omega_1+\omega_2+\cdots+\omega_n)}.
\end{equation*}
For $A = \{\text{exactly } k \text{ heads}\}$:
\begin{equation*}
  \PP(A) = p^k(1-p)^{n-k}\binom{n}{k}.
\end{equation*}
\end{example}

\begin{definition}[Probability space]
A probability space is a triple $(\Omega, \FF, \PP)$ where:
\begin{itemize}
  \item $\Omega$ is a non-empty set (sample space).
  \item $\FF$ is a \textbf{sigma-algebra} of subsets of $\Omega$:
    \begin{enumerate}[label=(\roman*)]
      \item $\Omega \in \FF$.
      \item If $A \in \FF$, then $A^c \in \FF$.
      \item If $A_1, A_2, \ldots, A_n \in \FF$, then $\bigcup_n A_n \in \FF$.
    \end{enumerate}
  \item $\PP : \FF \to [0,1]$ is a \textbf{probability measure} satisfying:
    \begin{equation*}
      \PP(\Omega) = 1, \qquad \PP\!\left(\bigsqcup_{i=1}^{\infty} A_i\right) = \sum_n \PP(A_n).
    \end{equation*}
\end{itemize}
\end{definition}

\begin{remark}
A discrete probability space can be subsumed into $(\Omega, \FF, \PP)$ by taking $\FF = 2^\Omega$ (the power set, i.e.\ the collection of all subsets of $\Omega$).
\end{remark}

\begin{definition}[Generated sigma-algebra]
Let $\mathcal{S}$ be a collection of subsets of $\Omega$. The \textbf{sigma-algebra generated by} $\mathcal{S}$ is
\begin{equation*}
  \sigma(\mathcal{S}) = \bigcap_{\substack{F \supseteq \mathcal{S} \\ F \text{ sigma-algebra}}} F,
\end{equation*}
i.e.\ the smallest sigma-algebra containing $\mathcal{S}$.
\end{definition}

\begin{definition}[Borel sigma-algebra]
For $\Omega = \RR$:
\begin{equation*}
  \mathcal{B}(\RR) = \sigma\bigl(\{(a,b) : a < b,\; a,b \in \RR\}\bigr).
\end{equation*}
More generally, $\mathcal{B}(\Omega) = \sigma(\{\text{open sets of } \Omega\})$.
\end{definition}

\begin{definition}[Random variable]
Given $(\Omega, \FF, \PP)$, a function $X : \Omega \to \RR$ is a \textbf{random variable} if for every $x \in \RR$,
\begin{equation*}
  \{\omega : X(\omega) \le x\} = X^{-1}((-\infty, x]) \in \FF.
\end{equation*}
\end{definition}

\begin{example}
Let $\Omega = \{1,2,\ldots,6\}$, $\FF = \{\emptyset, \{2,4,6\}, \{1,3,5\}, \Omega\}$, $X : \Omega \to \RR$ with
\begin{equation*}
  X(\omega) = \begin{cases} 1 & \text{if } \omega \bmod 3 = 0 \\ 0 & \text{otherwise.} \end{cases}
\end{equation*}
Then $X^{-1}((-\infty, 0]) = \{1,2,4,5\}$ and $\PP(\{1,2,4,5\}) = 0$.
\end{example}

\begin{definition}[Random process / Stochastic process]
A \textbf{random process} is a collection $X = (X_t : t \in T)$ of random variables.
\begin{itemize}
  \item $T$ finite $\Rightarrow$ random vectors.
  \item $T$ countable $\Rightarrow$ $\{0,1,2,3,\ldots\}$; $(X_1, X_2, \ldots)$ maps $(\Omega, \FF, \PP) \to (\RR^\infty, \mathcal{B}(\RR^\infty))$.
\end{itemize}
\end{definition}

\begin{definition}[Finite-dimensional distributions (fdd)]
For any $m \in \{1,2,3,\ldots\}$ and finite set $\{t_1, t_2, \ldots, t_m\} \subset T$, define
\begin{equation*}
  f_{t_1,t_2,\ldots,t_m}(x_1,x_2,\ldots,x_m) = \PP(X_{t_1} \le x_1,\, X_{t_2} \le x_2,\,\ldots,\, X_{t_m} \le x_m).
\end{equation*}
The collection $\{f_{t_1,\ldots,t_m} : m \in \{1,2,\ldots\},\, (t_1,\ldots,t_m)\}$ is called the \textbf{set of fdd}.

Example: $f_{(X,Y)}(x,y) = \PP(X \le x,\, Y \le y)$.
\end{definition}

\begin{definition}[Consistency condition]
If $T_m \subset T_n \subset T$, then
\begin{equation*}
  f_{T_m}(x_1,x_2,\ldots,x_m) = f_{T_n}(\infty,\ldots,\infty,x_1,\infty,\ldots,\infty),
\end{equation*}
where $\infty$ is placed at the coordinates in $T_n \setminus T_m$.
\end{definition}

\begin{theorem}[Kolmogorov Extension Theorem]
If the fdd are consistent, then they define a unique stochastic process.
\end{theorem}

\section{Discrete-Time Markov Chains}
\subsection{Definitions and Basic Properties}

\begin{definition}[Markov Chain (DTMC)]
A DTMC $X \sim \text{Markov}(\lambda, P)$ is a stochastic process $(X_n;\, n \ge 0)$ with:
\begin{enumerate}[label=(\roman*)]
  \item \textbf{State space} $\II$: a countable set.
  \item \textbf{Initial distribution} $\lambda = (\lambda_i : i \in \II)$ with $\sum_{i \in \II} \lambda_i = 1$, $0 \le \lambda_i \le \infty$, and $\PP(X_0 = i_0) = \lambda_{i_0}$.
  \item \textbf{Transition probability matrix} (stochastic matrix) $P = (p_{ij},\, i,j \in \II)$ with $\sum_j p_{ij} = 1$.
\end{enumerate}
The process satisfies: for all $n > 0$ and $i_1, i_2, \ldots, i_{n+1} \in \II$,
\begin{equation*}
  \PP(X_{n+1} = i_{n+1} \mid X_0 = i_0, X_1 = i_1, \ldots, X_n = i_n)
  = \PP(X_{n+1} = i_{n+1} \mid X_n = i_n) = p_{i_n i_{n+1}}.
\end{equation*}
\end{definition}

\begin{example}[Coin toss with bias $p$]
Let $X_n = 0$ (Tails) or $1$ (Heads). Then $X_0, X_1, \ldots, X_n, X_{n+1}$ are i.i.d.\ and the process is Markov.

Define $Y_0 = 0$ and for $n \ge 1$, $Y_n = \sum_{i=0}^{n-1} X_i$ (number of heads before the $n$-th toss). Check: $(Y_n)_{n \ge 0}$ is a DTMC.
\end{example}

\begin{theorem}[Characterisation via fdd]
$(X_n)_{n \ge 0} \sim \text{Markov}(\lambda, P)$ if and only if for all $N$ and $i_0, \ldots, i_N \in \II$,
\begin{equation*}
  \PP(X_0 = i_0, X_1 = i_1, \ldots, X_N = i_N) = \lambda_{i_0} p_{i_0 i_1} p_{i_1 i_2} \cdots p_{i_{N-1} i_N}.
\end{equation*}
\end{theorem}

\begin{proof}
The forward direction follows from the Markov property by induction:
\begin{align*}
  \PP(X_0 = i_0, \ldots, X_n = i_n) &= \PP(X_0 = i_0)\,\PP(X_1 = i_1 \mid X_0 = i_0)\cdots\PP(X_n = i_n \mid X_{n-1} = i_{n-1}) \\
  &= \lambda_{i_0} p_{i_0 i_1} \cdots p_{i_{n-1} i_n}.
\end{align*}
For the conditional probability:
\begin{align*}
  \PP(X_{n+1} = i_{n+1} \mid X_0 = i_0, \ldots, X_n = i_n)
  &= \frac{\PP(X_0 = i_0, \ldots, X_n = i_n, X_{n+1} = i_{n+1})}{\PP(X_0 = i_0, \ldots, X_n = i_n)} \\
  &= \frac{\lambda_{i_0} p_{i_0 i_1} \cdots p_{i_{n-1} i_n} \cdot p_{i_n i_{n+1}}}{\lambda_{i_0} p_{i_0 i_1} \cdots p_{i_{n-1} i_n}} = p_{i_n i_{n+1}}.
\end{align*}
\end{proof}

\begin{remark}
Write $\PP(\,\cdot \mid X_0 = i) = \PP_i(\,\cdot)$. Then $i \to j$ means $\PP_i(X_n = j \text{ for some } n \ge 0) > 0$.
\end{remark}

Define the Kronecker delta: $\delta_{ij} = \begin{cases}1 & i = j \\ 0 & i \neq j\end{cases}$ and $\delta_i = (\delta_{ij} : j \in \II)$.

\begin{theorem}[Markov property at time $m$]
Let $(X_n)_{n \ge 0} \sim \text{Markov}(\lambda, P)$. Then conditional on $X_m = i$, the process $(X_{m+n})_{n \ge 0} \sim \text{Markov}(\delta_i, P)$, and it is independent of $X_0, \ldots, X_m$.
\end{theorem}

\begin{proof}
We show: for all $A \in \sigma(X_0, X_1, \ldots, X_m)$,
\begin{equation*}
  \PP\bigl(\{X_m = i_m, \ldots, X_{m+n} = i_{m+n}\} \cap A \mid X_m = i\bigr)
  = \delta_{i,i_m} p_{i_m i_{m+1}} \cdots p_{i_{m+n-1} i_{m+n}} \,\PP(A \mid X_m = i).
\end{equation*}
\textbf{Step 1:} Take $A = \{X_0 = i_0, \ldots, X_m = i_m\}$. Using the fdd formula, the LHS gives:
\begin{align*}
  \frac{\lambda_{i_0} p_{i_0 i_1}\cdots p_{i_{m-1} i_m} \cdot p_{i_m i_{m+1}} \cdots p_{i_{m+n-1} i_{m+n}}}{\PP(X_m = i)}
  = \PP(A \cap \{X_m = i\}) \cdot \delta_{i,i_m} p_{i_m i_{m+1}} \cdots p_{i_{m+n-1} i_{m+n}} / \PP(X_m = i).
\end{align*}
\textbf{Step 2:} Extend to $A = \bigsqcup_{k=1}^\infty A_k$ by countable additivity. For $A = \{X_0 \in I_0, X_1 \in I_1, \ldots, X_m \in I_m\}$ where $I_0 = \{i_{01}, i_{02}\}$, we have $A = \{X_0 = i_{01}, \ldots, X_m = i_m\} \cup \{X_0 = i_{02}, \ldots, X_m = i_m\}$. The result follows by summing:
\begin{align*}
  &= \delta_{i,i_m} p_{i_m i_{m+1}} \cdots p_{i_{m+n-1} i_{m+n}} \sum_{k=1}^\infty \frac{\PP(A_k \cap \{X_m = i\})}{\PP(X_m = i)} = \PP\!\left(\bigsqcup_{k=1}^\infty A_k \cap \{X_m = i\}\right). \qquad \square
\end{align*}
\end{proof}

\begin{theorem}[Chapman--Kolmogorov Equations]
For any $n, l \ge 0$ and states $i, j \in \II$:
\[
  p_{ij}^{(n+l)} = \sum_{k \in \II} p_{ik}^{(n)}\, p_{kj}^{(l)}.
\]
In matrix form: $P^{n+l} = P^n \cdot P^l$.
\end{theorem}

\begin{proof}
By the law of total probability and the Markov property:
\begin{align*}
  p_{ij}^{(n+l)} &= \PP_i(X_{n+l} = j) \\
  &= \sum_{k \in \II} \PP_i(X_n = k)\, \PP_i(X_{n+l} = j \mid X_n = k) \\
  &= \sum_{k \in \II} p_{ik}^{(n)}\, \PP_k(X_l = j) \quad \text{(Markov property at time } n\text{)} \\
  &= \sum_{k \in \II} p_{ik}^{(n)}\, p_{kj}^{(l)}.
\end{align*}
The second equality uses the total probability law; the third uses the fact that given $X_n = k$, the future process $(X_{n+m})_{m\ge 0}$ is $\text{Markov}(\delta_k, P)$, independent of the past.
\end{proof}

\begin{remark}
The Chapman--Kolmogorov (CK) equation is the fundamental identity for computing $n$-step transition probabilities. It says: to go from $i$ to $j$ in $n+l$ steps, we must pass through some intermediate state $k$ after $n$ steps. The CK equations imply $P^n = P \cdot P^{n-1}$ (induction), so $p_{ij}^{(n)}$ can be computed as the $(i,j)$ entry of the matrix power $P^n$.
\end{remark}

\subsection{Class Structure}

The idea is to break down the Markov chain into smaller entities that can be analysed independently.

\begin{definition}[Accessibility / Leads to]
$i \to j$ ($i$ \emph{leads to} $j$, or $j$ is \emph{reachable} from $i$) if $p_{ij}^{(n)} > 0$ for some $n \ge 0$, i.e.\ $\PP_i(X_n = j \text{ for some } n \ge 0) > 0$.
\end{definition}

\begin{definition}[Communication]
$i \leftrightarrow j$ ($i$ \emph{communicates with} $j$) if $i \to j$ and $j \to i$.
\end{definition}

\begin{proposition}
$\leftrightarrow$ is an equivalence relation. Its equivalence classes are called \emph{communicating classes}.
\end{proposition}

\begin{remark}
$(P^2)_{ij} = \sum_{k \in \II} p_{ik} p_{kj}$.
\end{remark}

\begin{definition}[Closed communicating class]
A communicating class $C$ is \emph{closed} if for all $i \in C$ and $j \notin C$, $p_{ij} = 0$. (Once entered, cannot leave.)
\end{definition}

\begin{remark}[Diagram]
Example with 4 states: States $\{1,2,3,4\}$. States 1 and 2 communicate ($1 \leftrightarrow 2$), states 3 and 4 communicate ($3 \leftrightarrow 4$), with transitions from $\{1,2\}$ to $\{3,4\}$ possible. Class structure depends only on positive entries of $P$.
\end{remark}

\begin{definition}[Absorbing state]
A set of states $A$ is \emph{absorbing} if $A$ is a closed communicating class. If $A = \{i\}$, then $i$ is called an absorbing state (i.e.\ $p_{ii} = 1$).
\end{definition}

\begin{definition}[Irreducibility]
$P$ (or the Markov chain $(X_n)$) is said to be \emph{irreducible} if $\II$ is a single closed communicating class, i.e.\ $i \leftrightarrow j$ for all $i, j \in \II$.
\end{definition}

\subsection{Stopping Times and Strong Markov Property}

\begin{definition}[Stopping time]
A random time $T : \Omega \to \{0,1,2,\ldots\} \cup \{\infty\}$ is a \emph{stopping time} for $(X_n)_{n \ge 0} \sim \text{Markov}(\lambda, P)$ if for every $m \ge 0$,
\begin{equation*}
  \{T = m\} \in \sigma(X_0, X_1, \ldots, X_m),
\end{equation*}
i.e.\ the event $\{T = m\}$ depends only on $X_0, X_1, \ldots, X_m$.
\end{definition}

\begin{example}
Let $T = \min\{n \ge 0 : X_n = 20\}$ (first time to reach state 20). Then
\begin{equation*}
  \{T = m\} = \{X_0 \neq 20, X_1 \neq 20, \ldots, X_{m-1} \neq 20, X_m = 20\} \in \sigma(X_0, \ldots, X_m).
\end{equation*}
So $T$ is a stopping time.
\end{example}

\begin{theorem}[Strong Markov Property]
Let $(X_n)_{n \ge 0} \sim \text{Markov}(\lambda, P)$ and let $T$ be a stopping time. Then conditional on $T < \infty$ and $X_T = i$, the process $(X_{T+n})_{n \ge 0} \sim \text{Markov}(\delta_i, P)$, and it is independent of $X_0, \ldots, X_T$.
\end{theorem}


\begin{proof}
We must show that for any $n \ge 0$ and states $j_0, j_1, \ldots, j_n$,
\begin{equation*}
  \PP(X_{T+1}=j_1,\ldots,X_{T+n}=j_n \mid T < \infty, X_T = i, X_0,\ldots,X_{T-1})
  = P^n(i, j_n)\cdot\mathbf{1}[j_0 = i].
\end{equation*}
Condition on $\{T = m\}$ for each $m$. On $\{T = m\}$, the event $\{X_T = i\}$ is $\{X_m = i\}$, which is $\sigma(X_0,\ldots,X_m)$-measurable. By the ordinary Markov property applied at fixed time $m$, the future $(X_{m+1}, X_{m+2}, \ldots)$ is independent of $(X_0, \ldots, X_m)$ given $X_m = i$ and follows $\text{Markov}(\delta_i, P)$. Summing over all $m$ gives the result.
\end{proof}

\subsection{Hitting Times and Gambler's Ruin}

\begin{definition}[Hitting time / First passage time]
For a state $j \in \II$, define the \emph{first passage time} to $j$ as
\begin{equation*}
  T_j = \min\{n \ge 1 : X_n = j\}.
\end{equation*}
(If $X_n \neq j$ for all $n \ge 1$, set $T_j = \infty$.)

Define:
\begin{align*}
  h_j^i &= \PP_i(T_j < \infty) \quad \text{(probability of ever reaching } j \text{ starting from } i\text{)},\\
  m_j^i &= \EE_i[T_j] \quad \text{(expected hitting time of } j \text{ from } i\text{)}.
\end{align*}
\end{definition}

\begin{remark}
By convention, $\PP_i(\cdot)$ means probability given $X_0 = i$, i.e.\ with initial distribution $\lambda = \delta_i$.
\end{remark}

\subsubsection{First Step Analysis}

The key technique is to condition on $X_1$ (the first step) and use the Markov property. For the hitting probability $h_j^i$:
\begin{align*}
  h_j^i &= \PP_i(T_j < \infty) \\
         &= \sum_{k \in \II} p_{ik}\, \PP_i(T_j < \infty \mid X_1 = k)\\
         &= p_{ij} + \sum_{k \neq j} p_{ik}\, h_j^k.
\end{align*}
This gives the system of equations:
\begin{equation*}
  h_j^i = p_{ij} + \sum_{k \neq j} p_{ik}\, h_j^k, \quad i \neq j,
\end{equation*}
with boundary condition $h_j^j = 1$.

Similarly, for expected hitting times $m_j^i$:
\begin{equation*}
  m_j^i = 1 + \sum_{k \neq j} p_{ik}\, m_j^k, \quad i \neq j,
\end{equation*}
with $m_j^j = 0$.

\begin{example}[Gambler's Ruin]
A gambler starts with \$\(a\) ($0 < a < N$). At each step, with probability $p$ she wins \$1 and with probability $q = 1-p$ she loses \$1. The game stops when she reaches \$0 (ruin) or \$\(N\) (success). Model: $X_n$ = fortune at time $n$, $\II = \{0, 1, \ldots, N\}$, states 0 and $N$ are absorbing.

Let $h_i = \PP_i(\text{reach } N \text{ before } 0)$. By first step analysis:
\begin{equation*}
  h_i = p\, h_{i+1} + q\, h_{i-1}, \quad 1 \le i \le N-1,
\end{equation*}
with $h_0 = 0$ and $h_N = 1$.

\textbf{Case 1: $p = q = 1/2$.} The equation becomes $h_i = \frac{1}{2}h_{i+1} + \frac{1}{2}h_{i-1}$, i.e.\ $h_{i+1} - h_i = h_i - h_{i-1}$. So $h_i$ is linear: $h_i = \frac{i}{N}$.

\textbf{Case 2: $p \neq q$.} The characteristic equation of $ph_{i+1} - h_i + qh_{i-1} = 0$ has roots $r = 1$ and $r = q/p$. General solution: $h_i = A + B(q/p)^i$. Applying boundary conditions:
\begin{equation*}
  h_i = \frac{1 - (q/p)^i}{1 - (q/p)^N}.
\end{equation*}
\end{example}

\begin{theorem}[Minimality of Hitting Probabilities]
\label{thm:minimality}
Let $A \subseteq \II$ be a target set and define $h^A_i = \PP_i(T_A < \infty)$ where $T_A = \min\{n \ge 0 : X_n \in A\}$. Then $\mathbf{h}^A = (h^A_i)_{i \in \II}$ is the \emph{minimal non-negative solution} to:
\[
  x_i = \begin{cases} 1 & i \in A, \\ \displaystyle\sum_{j \in \II} p_{ij}\, x_j & i \notin A. \end{cases}
\]
That is: (a) $\mathbf{h}^A$ satisfies the system, and (b) if $\mathbf{x} \ge 0$ also satisfies the system, then $x_i \ge h^A_i$ for all $i$.
\end{theorem}

\begin{proof}
\textbf{(a) $h^A$ satisfies the system.} For $i \in A$: $h^A_i = \PP_i(T_A < \infty) = 1$ (already there). For $i \notin A$, condition on the first step:
\[
  h^A_i = \PP_i(T_A < \infty) = \sum_{j \in \II} p_{ij}\, \PP_j(T_A < \infty) = \sum_{j \in \II} p_{ij}\, h^A_j.
\]

\textbf{(b) Minimality.} Let $\mathbf{x} \ge 0$ satisfy the system. We show $x_i \ge h^A_i$ for all $i$ by induction on $n$: we prove $x_i \ge \PP_i(T_A \le n)$ for all $n \ge 0$.

\emph{Base case} $n=0$: $\PP_i(T_A \le 0) = \mathbf{1}_{i \in A} \le x_i$ (since $x_i = 1$ for $i \in A$).

\emph{Inductive step}: Assume $x_j \ge \PP_j(T_A \le n)$ for all $j$. For $i \notin A$:
\[
  x_i = \sum_{j} p_{ij} x_j \ge \sum_j p_{ij}\, \PP_j(T_A \le n) = \PP_i(T_A \le n+1).
\]
Taking $n \to \infty$: $x_i \ge \PP_i(T_A < \infty) = h^A_i$.
\end{proof}

\begin{remark}
The minimality condition distinguishes $h^A$ from other solutions. For example, on an irreducible recurrent chain, $x_i \equiv 1$ also satisfies the system (with $A = \{j\}$), and $h^A_i = 1$ is indeed the minimal solution. For transient chains, other solutions with $x_i > h^A_i$ may exist. The linear system alone does not pin down $h^A$; the minimality condition is essential.
\end{remark}

\subsection{Recurrence and Transience}

\begin{definition}[Recurrence and transience]
State $i$ is \emph{recurrent} if $\PP_i(X_n = i \text{ for infinitely many } n) = 1$, equivalently $h_i^i = 1$ (i.e.\ $T_i < \infty$ a.s.).

State $i$ is \emph{transient} if $\PP_i(X_n = i \text{ for infinitely many } n) = 0$, equivalently $h_i^i < 1$.
\end{definition}

\begin{theorem}[Recurrence criterion via return visits]
Let $V_i = \sum_{n=0}^{\infty} \mathbf{1}[X_n = i]$ be the total number of visits to $i$. Then:
\begin{enumerate}[label=(\roman*)]
  \item $i$ is recurrent $\iff$ $\EE_i[V_i] = \infty$ $\iff$ $\sum_{n=0}^{\infty} p_{ii}^{(n)} = \infty$.
  \item $i$ is transient $\iff$ $\EE_i[V_i] < \infty$ $\iff$ $\sum_{n=0}^{\infty} p_{ii}^{(n)} < \infty$.
\end{enumerate}
\end{theorem}

\begin{proof}
Note that $V_i = \sum_{n \ge 0} \mathbf{1}[X_n = i]$, so $\EE_i[V_i] = \sum_{n \ge 0} p_{ii}^{(n)}$.

Let $f = h_i^i = \PP_i(T_i < \infty)$. On each return to $i$, the chain restarts by the Strong Markov Property and returns again with probability $f$. So the number of visits $V_i$ has a geometric distribution: $\PP_i(V_i \ge k) = f^{k-1}$ for $k \ge 1$.

If $f = 1$ (recurrent): $\PP_i(V_i \ge k) = 1$ for all $k$, so $V_i = \infty$ a.s.

If $f < 1$ (transient): $\PP_i(V_i \ge k) = f^{k-1} \to 0$, so $V_i < \infty$ a.s.\ and $\EE_i[V_i] = 1/(1-f) < \infty$.
\end{proof}

\begin{theorem}[Class property of recurrence]
Recurrence and transience are class properties: if $i \leftrightarrow j$ then $i$ is recurrent iff $j$ is recurrent.
\end{theorem}

\begin{proof}
Suppose $i$ is recurrent and $i \leftrightarrow j$. Then there exist $m, n$ such that $p_{ij}^{(m)} > 0$ and $p_{ji}^{(n)} > 0$. For any $k \ge 0$:
\begin{equation*}
  p_{jj}^{(m+k+n)} \ge p_{ji}^{(n)}\, p_{ii}^{(k)}\, p_{ij}^{(m)}.
\end{equation*}
Summing over $k$: $\sum_k p_{jj}^{(m+k+n)} \ge p_{ji}^{(n)}\, p_{ij}^{(m)}\, \sum_k p_{ii}^{(k)} = \infty$, since $i$ is recurrent. Hence $j$ is recurrent.
\end{proof}

\begin{corollary}
In a finite irreducible Markov chain, all states are recurrent.
\end{corollary}

\begin{definition}[Positive recurrence and null recurrence]
A recurrent state $i$ is \emph{positive recurrent} if $m_i^i = \EE_i[T_i] < \infty$, and \emph{null recurrent} if $m_i^i = \infty$.
\end{definition}


\begin{theorem}[Recurrent class is closed]
Every recurrent communicating class is a closed class.
\end{theorem}

\begin{proof}
By contrapositive. Let $C$ be a communicating class that is not closed. Then $\exists\, i \in C$, $j \notin C$ with $p_{ij} > 0$. Since $j \notin C$, we have $j \not\to i$, so $\PP(T_i = \infty \mid X_0 = i) \ge p_{ij} > 0$. Hence $f_{ii} = \PP_i(T_i < \infty) < 1$, so $i$ is transient, contradicting recurrence of $C$.
\end{proof}

\begin{theorem}[Finite closed classes are recurrent]
If $C$ is a finite closed communicating class, then all states in $C$ are recurrent.
\end{theorem}

\begin{proof}
Suppose for contradiction some $i \in C$ is transient, so $f_{jj} < 1$ for all $j \in C$ (by class property). Let $V_j = \EE_i[V_j]$ where $V_j = \#$ visits to $j$. For any $i, j \in C$:
\begin{align*}
  \EE_i[V_j] &= \sum_{n=1}^{\infty} \PP_i(V_j \ge n) \\
             &\le 1 + \frac{f_{jj}}{1-f_{jj}} = \frac{1}{1-f_{jj}} < \infty.
\end{align*}
By the Strong Markov Property and summing over $j \in C$:
\begin{equation*}
  \sum_{j \in C} \EE_i[V_j] = \sum_{j \in C} \frac{1}{1-f_{jj}} < \infty.
\end{equation*}
But since $C$ is closed, the chain stays in $C$ forever, so $\sum_{j \in C} V_j = \infty$ a.s. Contradiction.
\end{proof}

\subsection{Stationary Distributions}

\begin{definition}[Stationary measure]
For a Markov chain $(X_n)_{n \ge 0}$ on $(\II, P)$, a row vector $\underline{\pi} = (\pi_i,\, i \in \II)$ with $\pi_i \ge 0$ is a \emph{stationary measure} if
\begin{equation*}
  \pi_j = \sum_{i \in \II} \pi_i\, p_{ij} \quad \forall\, j \in \II, \qquad \text{i.e.}\quad \underline{\pi} = \underline{\pi} P.
\end{equation*}
\end{definition}

\begin{definition}[Stationary / invariant distribution]
A stationary measure $\underline{\pi}$ is a \emph{stationary (invariant) distribution} if additionally $\sum_{i \in \II} \pi_i = 1$. Equivalently, $\underline{\pi} P^n = \underline{\pi}$ for all $n$.
\end{definition}

\begin{theorem}[Stationarity is preserved over time]
If $\underline{\pi}$ is a stationary distribution w.r.t.\ $P$ and $\PP(X_0 = i) = \pi_i$, then $\PP(X_m = i) = \pi_i$ for all $m$.
\end{theorem}

\begin{proof}
By induction. The base case $m=0$ holds by assumption. Assume $\PP(X_m = j) = \pi_j$ for all $j$. Then:
\begin{align*}
  \PP(X_{m+1} = i) &= \sum_{j \in \II} \PP(X_{m+1}=i \mid X_m=j)\,\PP(X_m=j)
  = \sum_{j \in \II} p_{ji}\,\pi_j = \pi_i.
\end{align*}
\end{proof}

\begin{example}[Two-state chain]
Let $\II = \{1,2\}$ with
\begin{equation*}
  P = \begin{pmatrix} 1-\alpha & \alpha \\ \beta & 1-\beta \end{pmatrix}, \quad \alpha,\beta \in (0,1).
\end{equation*}
The stationary distribution satisfies $\underline{\pi}P = \underline{\pi}$:
\begin{equation*}
  \underline{\pi} = \left(\frac{\beta}{\alpha+\beta},\, \frac{\alpha}{\alpha+\beta}\right).
\end{equation*}
\end{example}

\begin{theorem}[Limit of transition probabilities is stationary]
Let $\II$ be finite. Suppose for some $i \in \II$,
\begin{equation*}
  p_{ij}^{(n)} \to \pi_j \quad \text{as } n \to \infty, \quad \forall\, j \in \II.
\end{equation*}
Then $\underline{\pi} = (\pi_j : j \in \II)$ is a stationary distribution.
\end{theorem}

\begin{proof}
Since $\sum_{j} p_{ij}^{(n)} = 1$ and $|\II| < \infty$, taking $n\to\infty$ gives $\sum_j \pi_j = 1$. For stationarity:
\begin{align*}
  \sum_{j \in \II} \pi_j\, p_{jk} &= \lim_{n\to\infty} \sum_{j \in \II} p_{ij}^{(n)}\, p_{jk} = \lim_{n\to\infty} p_{ik}^{(n+1)} = \pi_k. \qedhere
\end{align*}
\end{proof}

\subsubsection{Stationary Measure via Cycle Formula}

\begin{definition}[Expected number of visits per cycle]
For a Markov chain starting at $k$, and any $i \in \II$, define
\begin{equation*}
  \rho_i^k = \EE_k\!\left[\sum_{n=0}^{T_k - 1} \mathbf{1}\{X_n = i\}\right],
\end{equation*}
i.e.\ the expected number of visits to state $i$ during one excursion (cycle) from $k$ back to $k$.
\end{definition}

\begin{theorem}[Cycle formula — stationary measure]
Let $P$ be irreducible and recurrent. Then:
\begin{enumerate}[label=(\alph*)]
  \item $\rho_k^k = 1$.
  \item $\underline{\rho}^k = (\rho_i^k : i \in \II)$ satisfies $\underline{\rho}^k P = \underline{\rho}^k$.
  \item $0 < \rho_i^k < \infty$ for all $i \in \II$.
\end{enumerate}
\end{theorem}

\begin{proof}
(a) Obvious: starting from $k$, we count the visit at time 0, so $\rho_k^k = 1$.

(b) We compute:
\begin{align*}
  \rho_i^k &= \EE_k\!\left[\sum_{n=1}^{T_k} \mathbf{1}\{X_n = i\}\right]
  = \sum_{n=1}^{\infty} \PP_k(X_n = i,\, T_k \ge n).
\end{align*}
Using the Markov property at time $n-1$:
\begin{align*}
  \rho_i^k &= \sum_{n=1}^{\infty}\sum_{j \in \II} p_{ji}\, \PP_k(X_{n-1}=j,\, T_k \ge n) \\
  &= \sum_{j \in \II} p_{ji}\sum_{m=0}^{\infty} \PP_k(X_m = j,\, T_k > m)
  = \sum_{j \in \II} \rho_j^k\, p_{ji}.
\end{align*}
Hence $\underline{\rho}^k P = \underline{\rho}^k$.

(c) Since $P$ is irreducible, $\exists\, n,m$ s.t.\ $p_{ik}^{(m)} > 0$ and $p_{ki}^{(n)} > 0$. From (b), $\underline{\rho}^k P^m = \underline{\rho}^k$, so $\rho_i^k \ge p_{ki}^{(n)}\cdot \rho_k^k = p_{ki}^{(n)} > 0$. Also $\sum_j \rho_j^k = \EE_k[T_k] = m_k < \infty$ if positive recurrent.
\end{proof}

\begin{lemma}[Uniqueness of stationary measure]
Let $P$ be irreducible and $\underline{\pi}$ an invariant measure with $\pi_k = 1$. Then $\underline{\pi} \ge \underline{\rho}^k$ (componentwise). If $P$ is recurrent then $\underline{\pi} = \underline{\rho}^k$.
\end{lemma}

\begin{definition}[Mean recurrence time]
For $i \in \II$, the \emph{mean recurrence time} is $m_i = \EE_i[T_i]$.
\end{definition}

\begin{definition}[Positive and null recurrence]
A state $j \in \II$ is:
\begin{itemize}
  \item \emph{positive recurrent} if $f_{jj} = 1$ and $m_j < \infty$,
  \item \emph{null recurrent} if $f_{jj} = 1$ and $m_j = \infty$.
\end{itemize}
\end{definition}

\begin{example}[Simple symmetric random walk on $\mathbb{Z}$]
The SSRW ($p = q = 1/2$) on $\mathbb{Z}$ is irreducible. It is known that all states are recurrent ($f_{jj}=1$) but null recurrent ($m_j = \infty$). Higher-dimensional symmetric random walks ($d \ge 3$) are transient.
\end{example}

\begin{theorem}[Expected Number of Visits Formula]
\label{thm:ENj}
Let $(X_n)_{n\ge 0} \sim \text{Markov}(\delta_i, P)$. Define $f_{ij} = \PP_i(T_j < \infty)$ and $f_{jj} = \PP_j(T_j < \infty)$. Let $N_j = \sum_{n=0}^\infty \mathbf{1}[X_n = j]$ be the total number of visits to $j$. Then:
\[
  \EE_i[N_j] = \frac{f_{ij}}{1 - f_{jj}}.
\]
In particular, for $i = j$: $\EE_j[N_j] = \frac{1}{1-f_{jj}}$ if $f_{jj} < 1$ (transient), and $\EE_j[N_j] = \infty$ if $f_{jj} = 1$ (recurrent).
\end{theorem}

\begin{proof}
Starting from $i$, the chain first reaches $j$ with probability $f_{ij}$. If it does reach $j$, then by the Strong Markov Property, from $j$ it makes subsequent returns to $j$, each with probability $f_{jj}$, independently.

Let $G = $ total visits to $j$ after the first visit. Then $G \sim \text{Geom}(1-f_{jj})$ (number of successes before first failure, with success probability $f_{jj}$), so $\EE[G] = f_{jj}/(1-f_{jj})$.

Therefore:
\begin{align*}
  \EE_i[N_j] &= \underbrace{f_{ij}}_{\PP(\text{ever visit }j)} \cdot \underbrace{\left(1 + \frac{f_{jj}}{1-f_{jj}}\right)}_{\EE[\text{visits to }j \mid \text{first visit made}]} \\[4pt]
  &= f_{ij} \cdot \frac{1}{1-f_{jj}}.
\end{align*}

More formally: the number of visits to $j$ is $N_j = \sum_{k=1}^\infty \mathbf{1}[V_k]$ where $V_k$ is the event of a $k$-th visit to $j$. By the Strong Markov Property, $\PP_i(V_k) = f_{ij} \cdot f_{jj}^{k-1}$ for $k \ge 1$. Hence:
\[
  \EE_i[N_j] = \sum_{k=1}^\infty \PP_i(V_k) = \sum_{k=1}^\infty f_{ij}\, f_{jj}^{k-1} = \frac{f_{ij}}{1 - f_{jj}}.
\]
\end{proof}

\begin{remark}
This formula also yields the recurrence/transience criterion: $\EE_i[N_j] = \infty \iff f_{jj} = 1$ (for $i$ from which $j$ is reachable, i.e., $f_{ij} > 0$). Combined with $\EE_i[N_j] = \sum_n p_{ij}^{(n)}$, this gives the divergence test: $j$ is recurrent iff $\sum_n p_{jj}^{(n)} = \infty$.
\end{remark}

\begin{example}[Null Recurrence Despite Negative Expected Drift]
Consider the chain on $\II = \{0, 1, 2, \ldots\}$ (birth-death) with $p_0 = 1$ (reflecting barrier at 0) and for $i \ge 1$:
\[
  p_{i,i+1} = p_i, \quad p_{i,i-1} = 1-p_i, \quad \text{with } p_i = \frac{1}{2} + \frac{c}{i} \text{ for some } c > 0.
\]
The mean drift at state $i$ is $\EE[X_{n+1} - X_n \mid X_n = i] = 2p_i - 1 = 2c/i > 0$ (positive drift!), yet for suitable $c$, this chain can be null recurrent.

More relevant to PS6 Q2: consider $p_i = 1/2 - c/i^2$, giving \emph{negative drift} $-2c/i^2$ yet the chain is still null recurrent. The key point is:

\textbf{Negative drift does not imply positive recurrence.} The distinction is:
\begin{itemize}
  \item \textbf{Transient:} $f_{jj} < 1$, equivalently $\sum_n p_{jj}^{(n)} < \infty$. The chain drifts to $+\infty$ and never returns.
  \item \textbf{Null recurrent:} $f_{jj} = 1$ but $m_j = \EE_j[T_j] = \infty$. The chain always returns, but expected return time is infinite. Stationary distribution does NOT exist.
  \item \textbf{Positive recurrent:} $f_{jj} = 1$ and $m_j < \infty$. Stationary distribution exists.
\end{itemize}
A sufficient condition for null recurrence (not positive recurrence): by Foster--Lyapunov, the chain is recurrent if $\EE[V(X_{n+1}) - V(X_n) \mid X_n = i] \le 0$ for $i$ large (with $V(i) \to \infty$), but not positive recurrent if no solution to $\EE[V(X_{n+1})] - V(i) \le -\epsilon$ exists.

The simple symmetric random walk ($p_i = 1/2$, zero drift) on $\mathbb{Z}$ is the canonical null recurrent example: recurrent by Polya's theorem, yet $m_0 = \infty$.
\end{example}

\begin{theorem}[Characterisation of positive recurrence]
\label{thm:pos-rec}
Let $P$ be irreducible. The following are equivalent:
\begin{enumerate}[label=(\roman*)]
  \item Every state is positive recurrent.
  \item Some state is positive recurrent.
  \item $P$ has an invariant distribution $\underline{\pi}$; moreover if \emph{(iii)} holds, $\pi_i = \dfrac{1}{m_i}$.
\end{enumerate}
\end{theorem}

\begin{proof}[Proof sketch]
(i)$\Rightarrow$(ii) is obvious. 

(ii)$\Rightarrow$(iii): Let $k$ be positive recurrent. Then $\underline{\rho}^k$ is a stationary measure with $\sum_j \rho_j^k = \EE_k[T_k] = m_k < \infty$. Normalizing gives $\underline{\pi} = \underline{\rho}^k / m_k$.

(iii)$\Rightarrow$(i): If $\underline{\pi}$ is an invariant distribution for irreducible $P$, then $P$ is recurrent (by Lemma above, $\underline{\pi} = \underline{\rho}^k/m_k$ implies $m_k < \infty$). By class property, all states are positive recurrent and $\pi_i = 1/m_i$.
\end{proof}


\subsection{Periodicity and Convergence}

\begin{definition}[Period of a state]
The \emph{period} of state $j \in \II$ is
\begin{equation*}
  d(j) = \gcd\{n \ge 1 : p_{jj}^{(n)} > 0\}.
\end{equation*}
A state $j$ is \emph{aperiodic} if $d(j) = 1$.
\end{definition}

\begin{example}
Consider the two-state chain with $P = \begin{pmatrix}0&1\\1&0\end{pmatrix}$. Then $P^2 = I$, $P^{2n} = I$, $P^{2n+1} = P$, so $d(1) = d(2) = 2$ (periodic with period 2).

A state $j$ is aperiodic iff $p_{jj}^{(n)} > 0$ for all sufficiently large $n$.
\end{example}

\begin{lemma}[Aperiodicity is a class property]
Suppose $P$ is irreducible and one state is aperiodic. Then for all $j, k \in \II$, $p_{jk}^{(n)} > 0$ for all sufficiently large $n$. In particular, all states are aperiodic.
\end{lemma}

\begin{remark}
Periodicity is a class property: if $i \leftrightarrow j$ then $d(i) = d(j)$.
\end{remark}

\begin{proof}[Proof sketch (aperiodicity lemma)]
Since $P$ is irreducible, $\exists\, r, s$ s.t.\ $p_{ji}^{(r)} > 0$ and $p_{ik}^{(s)} > 0$. Then
\begin{equation*}
  p_{jj}^{(r+n+s)} \ge p_{ji}^{(r)}\, p_{ii}^{(n)}\, p_{ik}^{(s)} > 0
\end{equation*}
for all sufficiently large $n$ (since $i$ is aperiodic). Hence $p_{jk}^{(r+n+s)} > 0$ for all large enough $n$.
\end{proof}

\begin{theorem}[Markov Chain Convergence Theorem]
\label{thm:convergence}
Let $P$ be irreducible and aperiodic, and suppose $\underline{\pi}$ is an invariant distribution for $P$. Then for any initial distribution $\lambda$,
\begin{equation*}
  \PP(X_n = j) \to \pi_j \quad \text{as } n \to \infty, \quad \forall\, j \in \II.
\end{equation*}
In particular, $p_{ij}^{(n)} \to \pi_j$ as $n \to \infty$, for all $i, j$.
\end{theorem}

\begin{proof}[Proof (Coupling argument)]
Define $(Y_n)_{n \ge 0} \sim \text{Markov}(\underline{\pi}, P)$ independently of $(X_n)$. Let $W_n = (X_n, Y_n)$; this is a Markov chain on $\II \times \II$ with transition probabilities
\begin{equation*}
  \tilde{P}_{(i,k)(j,\ell)} = p_{ij}\,p_{k\ell}.
\end{equation*}
The initial distribution of $W_0$ is $\mu_{(i,k)} = \lambda_i\,\pi_k$. Since $P$ is irreducible and aperiodic, $\tilde{P}$ is irreducible and positive recurrent with stationary distribution $\tilde{\pi}_{(i,k)} = \pi_i\,\pi_k$.

Fix $b \in \II$ and let $T = \inf\{n \ge 1 : X_n = Y_n = b\}$ (the \emph{coupling time}).

\textbf{Step 1:} $\PP(T < \infty) = 1$, since $\tilde{P}$ is positive recurrent (irreducible, finite or positive recurrent).

\textbf{Step 2:} Define the ``stitched'' process
\begin{equation*}
  Z_n = \begin{cases} X_n & \text{if } n < T, \\ Y_n & \text{if } n \ge T. \end{cases}
\end{equation*}
By the Strong Markov Property applied to $W_n$ at time $T$, $(Z_n)_{n \ge 0} \sim \text{Markov}(\lambda, P)$ and $(Z_n)$ has the same distribution as $(X_n)$.

\textbf{Step 3:} 
\begin{align*}
  |\PP(X_n = j) - \PP(Y_n = j)| &= |\PP(Z_n = j) - \PP(Y_n = j)| \\
  &= |\PP(X_n = j,\, n < T) - \PP(Y_n = j,\, n < T)| \\
  &\le \PP(n < T) \to 0 \quad \text{as } n \to \infty.
\end{align*}
Since $\PP(Y_n = j) = \pi_j$ for all $n$ (stationarity), $\PP(X_n = j) \to \pi_j$.
\end{proof}

\subsection{Foster--Lyapunov Criterion}

\begin{theorem}[Foster--Lyapunov Criterion]
\label{thm:foster-lyapunov}
Let $(X_n)$ be a Markov chain on $\II$ with transition matrix $P$.
\begin{enumerate}[label=(\roman*)]
  \item \textbf{Recurrence:} $P$ is recurrent iff $\exists$ a function $V : \II \to \RR_+$ and a finite set $A \subset \II$ such that:
  \begin{enumerate}[label=(\alph*)]
    \item $V(j) \to \infty$ as $j \to \infty$ (i.e.\ $V$ is a Lyapunov function),
    \item $\forall\, i \notin A$: $\EE[V(X_{m+1}) - V(X_m) \mid X_m = i] \le 0$ (negative drift outside $A$).
  \end{enumerate}
  \item \textbf{Transience:} $P$ is transient iff $\exists\, V : \II \to \RR_+$ and finite $A \subset \II$ such that:
  \begin{enumerate}[label=(\alph*)]
    \item $\exists\, j \in A$ s.t.\ $V(j) < \inf_{i \notin A} V(i)$,
    \item $\forall\, i \notin A$: $\EE[V(X_{m+1}) - V(X_m) \mid X_m = i] \le 0$.
  \end{enumerate}
\end{enumerate}

The \emph{mean drift} of $V$ from state $i$ in one step is:
\begin{equation*}
  \EE[V(X_{m+1}) - V(X_m) \mid X_m = i] = \sum_{j \in \II} p_{ij}(V(j) - V(i)).
\end{equation*}
\end{theorem}

\begin{example}[Random walk on $\NN$]
Consider a random walk on $\NN_0 = \{0, 1, 2, \ldots\}$ with $p_{i,i+1} = p$ and $p_{i,i-1} = q$ for $i \ge 1$, absorbing at 0. Take $V(n) = n$ and $A = \{0\}$. For $i \ge 1$:
\begin{equation*}
  \EE[V(X_{m+1}) - V(X_m) \mid X_m = i] = p(i+1) + q(i-1) - i = p - q.
\end{equation*}
If $p \le q$ (i.e.\ $p \le 1/2$), drift $\le 0$ for $i \notin A$. By (i), the chain is recurrent.

For transience, take $V(n) = \alpha^n$ with $0 < \alpha < 1$, $A = \{0\}$. Then:
\begin{align*}
  \EE[V(X_{m+1})-V(X_m)\mid X_m=i] &= p\alpha^{i+1} + q\alpha^{i-1} - \alpha^i = \alpha^{i-1}(p\alpha^2 - \alpha + q).
\end{align*}
This is $\le 0$ iff $p\alpha^2 - \alpha + q \le 0$. If $p > q$, choosing $\alpha = q/p < 1$ gives $p(q/p)^2 - q/p + q = q^2/p - q/p + q = q(q-1)/p + q \cdot(1 - 1/p + 1) $... the quadratic $p\alpha^2 - \alpha + q$ has roots $\alpha = 1$ and $\alpha = q/p$, so for $\alpha \in (q/p, 1)$ the expression is negative, confirming transience when $p > q$.
\end{example}


\subsection{Positive Recurrence via Foster--Lyapunov: Additional Cases}

Continuing the Foster--Lyapunov criterion: part (iii) of the theorem states that $P$ is \emph{positive recurrent} iff there exists $V:\mathcal{I}\to\mathbb{R}_+$, $\epsilon>0$, and a finite set $A\subset\mathcal{I}$ such that:
\begin{enumerate}[label=(\alph*)]
  \item $\forall i\in A$: $\EE[V(X_{m+1})\mid X_m=i]<\infty$.
  \item $\forall i\notin A$: $\EE[V(X_{m+1})-V(X_m)\mid X_m=i]\le -\epsilon$.
\end{enumerate}

\begin{example}[Positive recurrence: $V(x)=x$]
Take $V(x)=x$, $A=\{0\}$, drift $= p-q \le -\epsilon$ for $p<q$. So $p\le q \Rightarrow p<1/2$ implies positive recurrence.
\end{example}

\begin{example}[Foster--Lyapunov with $V(i)=i$, non-homogeneous drift]
Let $p_{0,1}=1$, $p_{i,i+1}=\tfrac{1}{2}-\tfrac{1}{2(i+1)}$ and $p_{i,i-1}=\tfrac{1}{2}+\tfrac{1}{2(i+1)}$ for $i\ge 1$. Take $V(i)=i$, $A=\{0\}$. Then:
\[
  \Delta V(i) = \EE[V(X_{m+1})-V(X_m)\mid X_m=i] = -\frac{1}{i+1}.
\]
This $\to 0$ as $i\to\infty$; the Foster--Lyapunov drift condition fails for a fixed $\epsilon>0$.

Take instead $V(i)=i^2$:
\[
  \EE[X_{m+1}^2 - X_m^2\mid X_m=i] = -\frac{i}{i+1} \le -\frac{1}{2} \quad\text{for all }i\ge 1.
\]
Foster--Lyapunov applies with $V(i)=i^2$ and $\epsilon=1/2$, so the chain is positive recurrent.
\end{example}

\subsection{Time Reversibility}

\begin{definition}[Time Reversal]
Let $(X_n)_{n\ge 0}\sim\text{Markov}(\pi,P)$ where $\pi$ is the stationary distribution. For any $N\ge 1$, the \emph{time-reversed chain} is $Y_n = X_{N-n}$, $0\le n\le N$. Then $(Y_n)_{0\le n\le N}$ is also a DTMC (by the Markov property applied in reverse).
\end{definition}

\begin{proof}[Derivation of reversed transition probabilities]
\begin{align*}
  \PP(Y_n = j\mid Y_{n+1} = i) &= \PP(X_{N-n}=j\mid X_{N-n-1}=i)\\
    &= \PP(X_{N-n-1}=i\mid X_{N-n}=j)\frac{\PP(X_{N-n}=j)}{\PP(X_{N-n-1}=i)}\\
    &= p_{ji}\frac{\PP(X_{N-n}=j)}{\PP(X_{N-n-1}=i)}.
\end{align*}
If the chain starts in stationarity ($\lambda=\pi$), then $\PP(X_n=j)=\pi_j$ for all $n$, so:
\[
  \PP(Y_n=j\mid Y_{n+1}=i) = \frac{p_{ji}\,\pi_j}{\pi_i}.
\]
The reversed chain has transition matrix $\hat{P}$ with $\hat{p}_{ij} = p_{ji}\pi_j/\pi_i$.
\end{proof}

\begin{definition}[Time Reversibility and Detailed Balance]
A DTMC with tpm $P$ and stationary distribution $\pi$ is \emph{time reversible} if the reversed chain has the same tpm as the original, i.e.,
\[
  \pi_i\,p_{ij} = \pi_j\,p_{ji} \quad \forall\,i,j\in\mathcal{I}.
\]
These are called the \textbf{Detailed Balance Equations (DBE)}.
\end{definition}

\begin{remark}
If the DBE hold, then summing over $j$:
\[
  \sum_{j\in\mathcal{I}}\pi_i p_{ij} = \sum_{j\in\mathcal{I}}\pi_j p_{ji} \implies \pi_i = \sum_{j}\pi_j p_{ji} \implies \pi P = \pi.
\]
So any $\pi$ satisfying the DBE is automatically a stationary distribution. This gives a much simpler route to finding $\pi$ in reversible chains: solve $\pi_i p_{ij}=\pi_j p_{ji}$.
\end{remark}

\begin{example}[Birth-death chain: advantages of reversibility]
Consider a birth-death chain on $\{0,1,2,3\}$ with $p_{i,i+1}=p$ and $p_{i,i-1}=q$ for $0<i<3$, reflecting at boundaries. The DBE give:
\[
  \pi_0 p = \pi_1 q,\quad \pi_1 p = \pi_2 q,\quad \pi_2 p = \pi_3 q.
\]
So $\pi_1 = \pi_0(p/q)$, $\pi_2 = \pi_0(p/q)^2$, $\pi_3 = \pi_0(p/q)^3$. Normalizing:
\[
  \pi_0\Bigl(1+(p/q)+(p/q)^2+(p/q)^3\Bigr) = 1.
\]
\textbf{Simple random walk on any graph is time reversible.}
\end{example}

\section{Borel--Cantelli Lemma and Ergodic Theorem}

\subsection{Borel--Cantelli Lemma}

\begin{definition}[Limsup and liminf of events]
For a sequence of events $\{A_n\}$:
\begin{align*}
  \limsup_{n} A_n &= \bigcap_{N\ge 1}\bigcup_{n\ge N} A_n = \{\omega: \omega\in A_n\text{ infinitely often}\} = \{A_n\text{ i.o.}\},\\
  \liminf_{n} A_n &= \bigcup_{N\ge 1}\bigcap_{n\ge N} A_n = \{\omega:\omega\notin A_n\text{ for only finitely many }n\}.
\end{align*}
\end{definition}

\begin{theorem}[Borel--Cantelli Lemma]
Let $\{A_n\}$ be a sequence of events.
\begin{enumerate}[label=(\roman*)]
  \item If $\sum_{n=1}^\infty \PP(A_n)<\infty$, then $\PP(A_n\text{ i.o.})=0$.
  \item If $\{A_n\}$ are independent and $\sum_{n=1}^\infty\PP(A_n)=\infty$, then $\PP(A_n\text{ i.o.})=1$.
\end{enumerate}
\end{theorem}

\begin{proof}[Proof of part (i)]
Let $B_N = \bigcup_{n\ge N}A_n$. Then $B_1\supseteq B_2\supseteq\cdots$ and $\bigcap_{N=1}^\infty B_N = \limsup_n A_n$. By continuity of probability from above:
\[
  \PP(A_n\text{ i.o.}) = \lim_{N\to\infty}\PP(B_N) = \lim_{N\to\infty}\PP\!\Bigl(\bigcup_{n\ge N}A_n\Bigr) \le \lim_{N\to\infty}\sum_{n=N}^\infty\PP(A_n) = 0,
\]
since $\sum_{n=1}^\infty\PP(A_n)<\infty$ implies the tail sums $\to 0$.
\end{proof}

\begin{remark}
Note: part (ii) requires independence. For DTMC: if $\sum_{n=1}^\infty p_{ii}^{(n)}<\infty$ (transient), then $\PP_i(X_n=i\text{ i.o.})=0$ by BCL(i). But when $\sum_{n}p_{ii}^{(n)}=\infty$, the events $\{X_n=i\}$ are \emph{not} independent (since $\PP_k(X_n=k, X_{n+1}=k)\ne\PP_k(X_n=k)\PP_k(X_{n+1}=k)$), so BCL(ii) does not directly apply.
\end{remark}

\subsection{Ergodic Theorem for DTMCs}

\begin{definition}[Empirical visit count]
For a DTMC $(X_n)_{n\ge 0}$, define:
\[
  V_i(n) = \sum_{k=0}^{n-1}\mathbf{1}\{X_k=i\} \quad\text{(number of visits to state }i\text{ before time }n).
\]
\end{definition}

\begin{theorem}[Ergodic Theorem for DTMCs]
Let $P$ be irreducible and let $\lambda$ be any initial distribution. If $(X_n)_{n\ge 0}\sim\text{Markov}(\lambda,P)$, then:
\begin{enumerate}[label=(\alph*)]
  \item $\displaystyle\PP\!\left[\frac{V_i(n)}{n}\to\frac{1}{m_i}\text{ as }n\to\infty\right]=1$, where $m_i=\EE_i[T_i]$ (with $1/\infty=0$).
  \item In the positive recurrent case, for any bounded $f:\mathcal{I}\to\mathbb{R}$:
  \[
    \PP\!\left(\frac{1}{n}\sum_{k=0}^{n-1}f(X_k)\to\bar{f}\text{ as }n\to\infty\right)=1,
  \]
  where $\bar{f}=\sum_{i\in\mathcal{I}}\pi_i f_i$ and $\pi$ is the unique stationary distribution.
\end{enumerate}
\end{theorem}

\begin{proof}
\textbf{Part (a) — Transient case:} If $P$ is transient, then $V_i:=\sum_{k=0}^\infty\mathbf{1}\{X_k=i\}<\infty$ a.s.\ (since $\EE_i[V_i]=1/(1-f_{ii})<\infty$, hence $V_i<\infty$ a.s.). Thus $V_i(n)\le V_i$ and $V_i(n)/n\le V_i/n\to 0$, so $1/m_i=0$ (since $m_i=\EE_i[T_i]=\infty$ for transient $i$).

\textbf{Part (a) — Recurrent case:} Fix state $i$ and let $T=T_i$. Since $P$ is recurrent, $\PP(T<\infty)=1$. Let $S_i^{(k)}$ denote the length of the $k$-th excursion to state $i$; by the Strong Markov Property, $S_i^{(1)},S_i^{(2)},\ldots$ are i.i.d.\ with mean $m_i=\EE_i[T_i]$.

The number of complete excursions by time $n$ satisfies $V_i(n)\to\infty$ as $n\to\infty$ (recurrence). Moreover, the sum of excursion lengths gives the squeeze:
\[
  \frac{S_i^{(1)}+\cdots+S_i^{(V_i(n)-1)}}{V_i(n)} \;\le\; \frac{n}{V_i(n)} \;\le\; \frac{S_i^{(1)}+\cdots+S_i^{(V_i(n))}}{V_i(n)}.
\]
By the Strong Law of Large Numbers for i.i.d.\ sequences, both sides converge a.s.\ to $m_i$. Hence $n/V_i(n)\to m_i$ a.s., i.e., $V_i(n)/n\to 1/m_i = \pi_i$ a.s.

\textbf{Part (b) — Positive recurrent case:} Since $P$ is positive recurrent, $\pi_i=1/m_i>0$ and $\sum_i\pi_i=1$. For any $J\subseteq\mathcal{I}$:
\begin{align*}
  \left|\frac{1}{n}\sum_{k=0}^{n-1}f(X_k)-\bar{f}\right| &= \left|\sum_{i\in\mathcal{I}}\left(\frac{V_i(n)}{n}-\pi_i\right)f_i\right|\\
  &\le \sum_{i\in J}\left|\frac{V_i(n)}{n}-\pi_i\right| + \sum_{i\notin J}\left(\frac{V_i(n)}{n}+\pi_i\right)\\
  &\le 2\sum_{i\in J}\left|\frac{V_i(n)}{n}-\pi_i\right| + 2\sum_{i\notin J}\pi_i.
\end{align*}
Given $\epsilon>0$, choose $J$ finite so that $\sum_{i\notin J}\pi_i<\epsilon/4$. Then, since $V_i(n)/n\to\pi_i$ a.s.\ for each $i$, there exists $N=N(\omega)$ such that for $n\ge N(\omega)$:
\[
  \sum_{i\in J}\left|\frac{V_i(n)}{n}-\pi_i\right|<\epsilon/4.
\]
Thus for $n\ge N(\omega)$, the whole expression is $<\epsilon$. Since $\epsilon$ was arbitrary, the convergence holds a.s.
\end{proof}


\section{Continuous-Time Markov Chains (CTMCs)}

\subsection{Introduction and Setup}

\begin{definition}[Continuous-time stochastic process]
A collection of random variables $(X_t)_{t\ge 0}$ taking values in a discrete set $\mathcal{S}$ (the state space) is a \emph{continuous-time random process}. We assume the sample paths $t\mapsto X_t(\omega)$ are right-continuous.

The \emph{law} of $(X_t)$ is determined by its finite-dimensional distributions:
\[
  \PP\bigl(\{X_{t_1}=i_1,\,X_{t_2}=i_2,\ldots,X_{t_n}=i_n\}\bigr)
\]
for all $t_1<t_2<\cdots<t_n$ and states $i_1,\ldots,i_n$.
\end{definition}

\begin{definition}[Continuous-Time Markov Chain]
$(X_t)_{t\ge 0}$ is a \emph{Continuous-Time Markov Chain (CTMC)} if for all $t>0$:
\[
  (X_{t+s},\,s\ge 0)\;\perp\!\!\!\perp\;(X_r,\,r<t)\;\big|\;X_t.
\]
Equivalently, for $t_1<t_2<\cdots<t_n<t<s$:
\[
  \PP(X_s=j\mid X_t=i,\,X_{t_n}=i_n,\ldots,X_{t_1}=i_1) = \PP(X_s=j\mid X_t=i).
\]
\end{definition}

\begin{example}[DTMC embedded in continuous time]
Any DTMC $(X_n)$ with $p_{ii}>0$ for all $i$ can be embedded into a CTMC by holding at each state for an $\Exp(\lambda)$ random time.
\end{example}

\subsection{Exponential Distribution and Memoryless Property}

\begin{definition}[Exponential random variable]
A random variable $T:\Omega\to[0,\infty)$ has the \emph{exponential distribution} with rate $\lambda>0$, written $T\sim\Exp(\lambda)$, if $f_T(t)=\lambda e^{-\lambda t}$ for $t\ge 0$. Equivalently, $\PP(T>t)=e^{-\lambda t}$.
\end{definition}

The key property of the exponential distribution is the \emph{memoryless property}:
\[
  \PP(T>t+s\mid T>t) = \PP(T>s) \quad \forall\,t,s\ge 0. \tag{$*$}
\]

\begin{lemma}[Characterization of Exponential]
$T$ satisfies $(*)$ if and only if $T\sim\Exp(\lambda)$ for some $\lambda>0$.
\end{lemma}

\begin{proof}
$(\Rightarrow)$ is clear. $(\Leftarrow)$: Let $g(t)=\PP(T>t)$. Then $(*)$ gives $g(s+t)=g(s)\cdot g(t)$ for all $s,t\ge 0$. Since $T>0$, $g(1/n)>0$ for large enough $n$, and $g(1)=\bigl(g(1/n)\bigr)^n>0$. Let $g(1)=e^{-\lambda}$. Then $g(p/q)=\bigl(g(1/q)\bigr)^p=e^{-\lambda p/q}$ for all $p,q\in\mathbb{Z}$. Since $g$ is decreasing and right-continuous, $g(t)=e^{-\lambda t}$ for all real $t\ge 0$.
\end{proof}

\begin{example}[Minimum of exponentials]
If $S_1\sim\Exp(\lambda_1),\ldots,S_n\sim\Exp(\lambda_n)$ are independent, then
\[
  \min\{S_1,S_2,\ldots,S_n\} \sim \Exp(\lambda_1+\lambda_2+\cdots+\lambda_n).
\]
\end{example}

\subsection{Poisson Process}

\begin{definition}[Poisson Process]
A \emph{Poisson process} of rate $\lambda>0$, written $\text{PP}(\lambda)$, is a right-continuous increasing integer-valued process $(X_t)_{t\ge 0}$ with $X_0=0$, constructed as follows. Let $\{S_n\}_{n\ge 1}$ be i.i.d.\ $\Exp(\lambda)$ inter-arrival times. Set $J_n = S_1+\cdots+S_n$ (time of $n$-th jump). Then:
\[
  X_t = n \quad\text{when } J_n\le t < J_{n+1}.
\]
\end{definition}

\begin{theorem}[Markov Property of Poisson Process]
Let $(X_t)$ be a $\text{PP}(\lambda)$. For any $s>0$, $(X_{t+s}-X_s)_{t\ge 0}$ is a $\text{PP}(\lambda)$ and is independent of the past $(X_r)_{r\le s}$.
\end{theorem}

\begin{proof}
It suffices to show this conditioned on $\{X_s=i\}$. On $\{X_s=i\}$:
\begin{enumerate}[label=(\roman*)]
  \item $X_r = \sum_{k=1}^{i}\mathbf{1}\{J_k\le r\}$ for $r<s$.
  \item The shifted holding times $\tilde{S}_1 = S_{i+1}-(s-J_i)$, $\tilde{S}_2=S_{i+2}$, etc.\ are i.i.d.\ $\Exp(\lambda)$ by the memoryless property, and independent of $\{S_1,\ldots,S_i\}$.
\end{enumerate}
Hence $(X_{t+s}-X_s)_{t\ge 0}$ has i.i.d.\ $\Exp(\lambda)$ inter-arrival times, making it a $\text{PP}(\lambda)$ independent of $(X_r)_{r\le s}$.
\end{proof}

\begin{theorem}[Strong Markov Property of Poisson Process]
Let $T$ be a stopping time for $(X_t)$. Conditioned on $\{T<\infty\}$, the process $(X_{T+t}-X_T)_{t\ge 0}$ is a $\text{PP}(\lambda)$ independent of $(X_s,\,s\le T)$.
\end{theorem}

\begin{definition}[Stationary and Independent Increments]
A process $(X_t)$ has \emph{stationary increments} if the law of $(X_{t+s}-X_s)$ depends only on $t$ (not on $s$). It has \emph{independent increments} if increments over disjoint intervals are independent.
\end{definition}

\begin{theorem}[Characterization of Poisson Process]
Let $(X_t)$ be an increasing integer-valued process with $X_0=0$. The following are equivalent:
\begin{enumerate}[label=(\roman*)]
  \item The inter-arrival times $\{S_n\}$ are i.i.d.\ $\Exp(\lambda)$.
  \item $(X_t)$ has independent increments and, uniformly in $t$:
  \begin{align*}
    \PP(X_{t+h}-X_t=0) &= 1-\lambda h+o(h),\\
    \PP(X_{t+h}-X_t=1) &= \lambda h+o(h),
  \end{align*}
  as $h\to 0$.
  \item $(X_t)$ has independent and stationary increments and $X_t\overset{d}{=}\text{Poi}(\lambda t)$.
\end{enumerate}
\end{theorem}

\begin{proof}[Proof: $(i)\Rightarrow(ii)$]
By the Markov property, $(X_t)$ has independent increments. For the increment probabilities, note:
\begin{align*}
  \PP(X_{t+h}-X_t\ge 1) &= \PP(S_1\le h) = 1-e^{-\lambda h} = \lambda h+o(h),\\
  \PP(X_{t+h}-X_t\ge 2) &\le \PP(S_1<h,S_2<h) = (1-e^{-\lambda h})^2 = o(h).
\end{align*}
So $\PP(X_{t+h}-X_t=1)=\lambda h+o(h)$ and $\PP(X_{t+h}-X_t=0)=1-\lambda h+o(h)$.

\medskip\noindent\textbf{Proof: $(ii)\Rightarrow(iii)$} Let $p_j(t)=\PP(X_t=j)$ for $j\ge 0$. By independent increments:
\begin{align*}
  p_j(t+h) &= \sum_{i=0}^{j} p_{j-i}(t)\,\PP(X_{t+h}-X_t=i)\\
  &= p_j(t)(1-\lambda h+o(h)) + p_{j-1}(t)(\lambda h+o(h)) + o(h).
\end{align*}
This gives:
\[
  \frac{p_j(t+h)-p_j(t)}{h} = -\lambda p_j(t) + \lambda p_{j-1}(t) + \frac{o(h)}{h}.
\]
Letting $h\to 0$: $p_j'(t) = -\lambda p_j(t)+\lambda p_{j-1}(t)$, with $p_0(0)=1$, $p_j(0)=0$ for $j\ge 1$.

For $j=0$: $p_0'(t)=-\lambda p_0(t)$, so $p_0(t)=e^{-\lambda t}$.
For $j=1$: $p_1'(t)=-\lambda p_1(t)+\lambda e^{-\lambda t}$, so $p_1(t)=(\lambda t)e^{-\lambda t}$.
By induction: $p_j(t) = \dfrac{(\lambda t)^j e^{-\lambda t}}{j!}$, i.e., $X_t\sim\text{Poi}(\lambda t)$.

Since $(X_t)$ satisfies (ii) with increments depending only on interval length, it also has stationary increments.

\medskip\noindent\textbf{Proof: $(iii)\Rightarrow(i)$} Property (iii) determines all finite-dimensional distributions, and we know that a process satisfying (i) also satisfies (iii). Since (iii) determines the law, (i) follows.
\end{proof}


\begin{theorem}[Superposition of Poisson Processes]
If $N_1(t) \sim \text{PP}(\lambda_1)$ and $N_2(t) \sim \text{PP}(\lambda_2)$ are independent, then their superposition (merger) $N(t) = N_1(t) + N_2(t) \sim \text{PP}(\lambda_1 + \lambda_2)$.

More generally, the superposition of $k$ independent Poisson processes with rates $\lambda_1, \ldots, \lambda_k$ is $\text{PP}(\lambda_1 + \cdots + \lambda_k)$.
\end{theorem}

\begin{proof}
$N(t) = N_1(t) + N_2(t)$ has independent increments (inherited from $N_1, N_2$ independent). For the increment distribution:
\[
  N(t) - N(s) = (N_1(t)-N_1(s)) + (N_2(t)-N_2(s)),
\]
the sum of independent $\text{Poi}(\lambda_1(t-s))$ and $\text{Poi}(\lambda_2(t-s))$ random variables, which is $\text{Poi}((\lambda_1+\lambda_2)(t-s))$. By the characterization theorem, $N \sim \text{PP}(\lambda_1+\lambda_2)$.
\end{proof}

\begin{theorem}[Splitting (Thinning) of Poisson Processes]
Let $N(t) \sim \text{PP}(\lambda)$. Independently mark each arrival as type 1 with probability $p$ and type 2 with probability $1-p$. Let $N_1(t)$ and $N_2(t)$ count type-1 and type-2 arrivals respectively. Then:
\begin{enumerate}[label=(\roman*)]
  \item $N_1(t) \sim \text{PP}(\lambda p)$ and $N_2(t) \sim \text{PP}(\lambda(1-p))$.
  \item $N_1$ and $N_2$ are \textbf{independent} Poisson processes.
\end{enumerate}
\end{theorem}

\begin{proof}
We verify the joint distribution. For disjoint time intervals and any $m, n \ge 0$:
\begin{align*}
  \PP(N_1(t) = m,\; N_2(t) = n) &= \PP(N(t) = m+n) \cdot \binom{m+n}{m} p^m (1-p)^n \\
  &= \frac{e^{-\lambda t}(\lambda t)^{m+n}}{(m+n)!} \cdot \frac{(m+n)!}{m!\, n!} p^m (1-p)^n \\
  &= \frac{e^{-\lambda p t}(\lambda p t)^m}{m!} \cdot \frac{e^{-\lambda(1-p)t}(\lambda(1-p)t)^n}{n!}.
\end{align*}
This factors as $\PP(N_1(t)=m) \cdot \PP(N_2(t)=n)$, confirming independence and the marginal Poisson distributions. The same argument applies to increments over any interval, establishing independence of the processes (not just at fixed time $t$).
\end{proof}

\begin{example}[PS9 Q2 type: Splitting into independent streams]
A router receives packets as a $\text{PP}(\lambda)$. Each packet is independently routed to server A with probability $p$ and server B with probability $1-p$. By the splitting theorem, arrivals to server A form a $\text{PP}(\lambda p)$ and to server B form a $\text{PP}(\lambda(1-p))$, and these two streams are \emph{independent}. This independence is surprising (one might expect them to be negatively correlated), and is a direct consequence of the Poisson distribution.
\end{example}

\subsection{Generator Matrix and Transition Rate Matrix}

\begin{definition}[Generator (Infinitesimal Generator) of a CTMC]
Let $(X_t)_{t\ge 0}$ be a CTMC on state space $S$. The \textbf{generator matrix} (or \textbf{Q-matrix}) $Q=(q_{ij})_{i,j\in S}$ is defined by:
\begin{align*}
  q_{ij} &\ge 0 \quad \text{for } i\ne j \quad (\text{transition rate from } i \text{ to } j),\\
  q_{ii} &= -\sum_{j\ne i} q_{ij} \quad (\text{diagonal entries are negative}).
\end{align*}
The \textbf{total departure rate} from state $i$ is $\lambda_i = -q_{ii} = \sum_{j\ne i} q_{ij}$.
\end{definition}

\begin{remark}
The $Q$-matrix is the continuous-time analogue of the one-step transition matrix $P$ for DTMCs. Informally, for small $h>0$:
\[
  \PP(X_{t+h}=j \mid X_t=i) \approx \begin{cases} q_{ij}\,h & j\ne i,\\ 1 + q_{ii}\,h & j=i. \end{cases}
\]
Thus $P(h) \approx I + Qh$ for small $h$, motivating the name ``infinitesimal generator.''
\end{remark}

\begin{definition}[Transition Probability Matrix]
For a CTMC with generator $Q$, the \textbf{transition probability matrix} $P(t) = (p_{ij}(t))$ is defined by
\[
  p_{ij}(t) = \PP(X_{t+s}=j \mid X_s=i), \quad t\ge 0, \; i,j\in S.
\]
By time-homogeneity, $P(t)$ does not depend on $s$.
\end{definition}

\subsection{Kolmogorov Equations}

\begin{theorem}[Kolmogorov Backward Equation]
For a CTMC with generator $Q$ and transition matrices $P(t)$:
\[
  \frac{d}{dt} P(t) = Q\,P(t), \quad P(0) = I.
\]
\end{theorem}

\begin{theorem}[Kolmogorov Forward Equation]
Under appropriate regularity conditions:
\[
  \frac{d}{dt} P(t) = P(t)\,Q, \quad P(0) = I.
\]
\end{theorem}

\begin{proof}[Proof sketch (Backward Equation)]
Fix $t>0$ and let $h\to 0$. By the Chapman–Kolmogorov equation:
\[
  P(t+h) = P(h)\,P(t).
\]
Then:
\begin{align*}
  \frac{P(t+h)-P(t)}{h} &= \frac{P(h)-I}{h}\,P(t).
\end{align*}
Using $P(h) = I + Qh + o(h)$, we get $\dfrac{P(h)-I}{h} \to Q$ as $h\to 0$. Hence:
\[
  P'(t) = Q\,P(t).
\]
\end{proof}

\begin{theorem}[Matrix Exponential Solution]
The unique solution to both Kolmogorov equations is:
\[
  P(t) = e^{Qt} := \sum_{n=0}^{\infty} \frac{(Qt)^n}{n!}.
\]
\end{theorem}

\begin{remark}
The matrix exponential $e^{Qt}$ has the following key properties:
\begin{itemize}
  \item $e^{Q\cdot 0} = I$ (initial condition satisfied).
  \item All entries of $e^{Qt}$ are non-negative (since $Q$-matrix yields valid probabilities).
  \item Each row of $e^{Qt}$ sums to 1 (stochastic matrix).
  \item $\frac{d}{dt}e^{Qt} = Q\,e^{Qt} = e^{Qt}\,Q$ (both equations satisfied simultaneously).
\end{itemize}
\end{remark}

\subsection{Stationary Distribution of CTMCs}

\begin{definition}[Stationary Distribution for CTMC]
A probability distribution $\pi = (\pi_i)_{i\in S}$ is a \textbf{stationary distribution} for the CTMC with generator $Q$ if:
\[
  \pi Q = \mathbf{0} \quad \text{and} \quad \sum_{i\in S} \pi_i = 1.
\]
Equivalently, in component form: for each $j\in S$,
\[
  \sum_{i\in S} \pi_i\,q_{ij} = 0, \quad \text{i.e.,} \quad \pi_j\,\lambda_j = \sum_{i\ne j} \pi_i\,q_{ij}.
\]
This is the \textbf{global balance equation}: rate of leaving state $j$ = rate of entering state $j$.
\end{definition}

\begin{remark}[Relation to DTMC Stationary Distribution]
For DTMCs, stationarity was $\pi P = \pi$. For CTMCs, stationarity is $\pi Q = 0$. If $\pi$ is stationary for the CTMC, then for any $t\ge 0$:
\[
  \pi\,P(t) = \pi\,e^{Qt} = \pi\left(\sum_{n=0}^\infty \frac{(Qt)^n}{n!}\right) = \pi + \pi Q\,t + \cdots = \pi.
\]
since $\pi Q = 0$ implies all higher-order terms vanish too.
\end{remark}

\subsection{Birth--Death Processes}

\begin{definition}[Birth--Death Process]
A CTMC on $S = \{0,1,2,\ldots\}$ is a \textbf{birth--death process} if transitions only occur between adjacent states: for $i\ge 0$,
\[
  q_{i,i+1} = \lambda_i \ge 0 \quad (\text{birth rate}), \quad q_{i,i-1} = \mu_i \ge 0 \quad (\text{death rate}), \quad q_{i,i} = -(\lambda_i+\mu_i).
\]
All other off-diagonal entries are zero ($q_0$ has no death: $\mu_0=0$).
\end{definition}

\begin{theorem}[Stationary Distribution of Birth--Death Process]
For an irreducible birth--death process, the stationary distribution (if it exists) satisfies the \textbf{detailed balance equations}:
\[
  \pi_i\,\lambda_i = \pi_{i+1}\,\mu_{i+1}, \quad i=0,1,2,\ldots
\]
The solution is:
\[
  \pi_n = \pi_0 \prod_{k=0}^{n-1}\frac{\lambda_k}{\mu_{k+1}}, \quad n\ge 1,
\]
where $\pi_0$ is determined by normalization $\sum_{n=0}^\infty \pi_n = 1$.
\end{theorem}

\begin{proof}
The global balance equation for state $j$ in a birth--death chain reduces to:
\[
  (\lambda_j+\mu_j)\pi_j = \lambda_{j-1}\pi_{j-1} + \mu_{j+1}\pi_{j+1}.
\]
For birth--death processes, this telescopes: summing the equation over states $\{0,1,\ldots,j\}$ yields $\pi_j\lambda_j = \pi_{j+1}\mu_{j+1}$ (detailed balance). Hence:
\[
  \pi_{n} = \frac{\lambda_{n-1}}{\mu_n}\pi_{n-1} = \frac{\lambda_{n-1}\lambda_{n-2}\cdots\lambda_0}{\mu_n\mu_{n-1}\cdots\mu_1}\pi_0.
\]
\end{proof}

\subsection{The M/M/1 Queue}

\begin{example}[M/M/1 Queue]
The \textbf{M/M/1 queue} models a single-server queueing system with:
\begin{itemize}
  \item Poisson arrivals at rate $\lambda$ (Markovian inter-arrival times $\sim\Exp(\lambda)$),
  \item Exponential service times $\sim\Exp(\mu)$ (one server),
  \item Unlimited waiting room (infinite buffer), FCFS discipline.
\end{itemize}
The state $X_t = $ number of customers in the system at time $t$ is a birth--death CTMC with:
\[
  \lambda_n = \lambda \quad (n\ge 0), \qquad \mu_n = \mu \quad (n\ge 1), \qquad \mu_0 = 0.
\]
The \textbf{traffic intensity} is $\rho = \lambda/\mu$.
\end{example}

\begin{theorem}[Stationary Distribution of M/M/1]
The M/M/1 queue has a stationary distribution if and only if $\rho = \lambda/\mu < 1$. In that case:
\[
  \pi_n = (1-\rho)\,\rho^n, \quad n=0,1,2,\ldots
\]
(geometric distribution with parameter $\rho$).
\end{theorem}

\begin{proof}
Apply the detailed balance equations: $\pi_n\,\lambda = \pi_{n+1}\,\mu$, so $\pi_{n+1} = (\lambda/\mu)\pi_n = \rho\,\pi_n$.
By induction: $\pi_n = \rho^n\,\pi_0$.
Normalization: $\sum_{n=0}^\infty \rho^n\pi_0 = \pi_0/(1-\rho) = 1$ requires $\rho<1$, giving $\pi_0=1-\rho$.
Hence $\pi_n = (1-\rho)\rho^n$.
\end{proof}

\begin{remark}[Performance Measures for M/M/1]
Under stationarity ($\rho<1$), Little's Law ($L=\lambda W$) and the geometric stationary distribution give:
\begin{align*}
  L &= \EE[X] = \sum_{n=0}^\infty n\pi_n = \frac{\rho}{1-\rho} \quad \text{(mean number in system)},\\
  W &= L/\lambda = \frac{1}{\mu-\lambda} \quad \text{(mean sojourn time)},\\
  L_q &= \frac{\rho^2}{1-\rho} \quad \text{(mean queue length, excluding in service)},\\
  W_q &= \frac{\rho}{\mu-\lambda} \quad \text{(mean waiting time in queue)}.
\end{align*}
\end{remark}

\begin{example}[Poisson Process as Pure Birth Chain]
The Poisson process $N_t \sim \text{Poi}(\lambda t)$ is the simplest birth--death chain with $\lambda_n = \lambda$ for all $n$ and $\mu_n = 0$. Its $Q$-matrix is:
\[
Q = \begin{pmatrix} -\lambda & \lambda & 0 & 0 & \cdots \\ 0 & -\lambda & \lambda & 0 & \cdots \\ 0 & 0 & -\lambda & \lambda & \cdots \\ \vdots & & & \ddots \end{pmatrix}.
\]
\end{example}

\begin{example}[Tandem Queue]
A \textbf{tandem queue} consists of two servers in series with a single Poisson arrival stream at rate $\lambda$. Server 1 has $\Exp(\mu_1)$ service times; departures from server 1 feed server 2 with $\Exp(\mu_2)$ service times. The joint queue length $(X_t^1, X_t^2)$ is a CTMC on $\mathbb{Z}_+^2$.
\end{example}

\subsection{Explosion and Minimal Process}

\begin{definition}[Explosion Time]
For a CTMC $(X_t)$ with holding times $S_n \sim \Exp(\lambda_{Y_n})$ in state $Y_n$, the \textbf{explosion time} is:
\[
\zeta = \lim_{n\to\infty} J_n = \sum_{n=0}^\infty S_n,
\]
where $J_n = S_0 + S_1 + \cdots + S_{n-1}$ is the time of the $n$-th jump. If $\zeta < \infty$ with positive probability, the chain \textbf{explodes}. After explosion, we set $X_t = \infty$ (cemetery state), giving the \textbf{minimal process}.
\end{definition}

\begin{remark}
The ``minimal'' process is the canonical construction: it runs until explosion, then stays at $\infty$. For a CTMC $(X_t)$ that is strong Markov, the event $\{T \le t\}$ depends only on $(X_s, s \le t)$.
\end{remark}

\subsection{Class Structure of CTMCs}

\begin{definition}[Accessibility and Communication in CTMCs]
State $j$ is \textbf{accessible} from $i$ (written $i \to j$) if $\PP_i(X_t = j \text{ for some } t) > 0$.
States $i$ and $j$ \textbf{communicate} ($i \leftrightarrow j$) if $i \to j$ and $j \to i$.
Communicating classes, irreducibility, and closed sets are defined as in DTMCs.
\end{definition}

\begin{theorem}[Class Structure Equivalences]
For states $i, j$ in a CTMC, the following are equivalent:
\begin{enumerate}[label=(\roman*)]
  \item $i \to j$ in $(X_t)$.
  \item $i \to j$ in the embedded DTMC $\{Y_n\}$.
  \item $q_{i,i_1} q_{i_1,i_2} \cdots q_{i_{n-1},j} > 0$ for some states $i_1,\ldots,i_{n-1} \in S$.
  \item $p_{ij}(t) > 0$ for all $t > 0$.
  \item $p_{ij}(t) > 0$ for some $t > 0$.
\end{enumerate}
\end{theorem}

\begin{proof}
$(iv) \Rightarrow (v) \Rightarrow (i)$ are clear. $(i) \Rightarrow (ii)$ is an exercise.

$(ii) \Rightarrow (iii)$: There exist states $i_1,\ldots,i_{n-1}$ such that $p_{i,i_1}p_{i_1,i_2}\cdots p_{i_{n-1},j} > 0$, which is equivalent to $\frac{q_{i,i_1}}{q_i}\cdots\frac{q_{i_{n-1},j}}{q_{i_{n-1}}} > 0$.

$(iii) \Rightarrow (iv)$: Suppose $q_{ij} > 0$. By the Chapman--Kolmogorov equation:
\[
p_{ij}(t) \ge \PP_i(J_1 < t,\; X_{J_1} = j,\; S_2 > t) = \PP_i(J_1 < t,\; X_{J_1} = j)\,\PP_i(S_2 > t) > 0.
\]
For a general chain of length $n$, iterate using Chapman--Kolmogorov: $p_{ij}(t) \ge p_{i,i_1}(t/n)\cdots p_{i_{n-1},j}(t/n) > 0$ for all $t > 0$.
\end{proof}

\begin{remark}
Condition (iv) shows that CTMCs have \textbf{no notion of periodicity}: once $i \to j$, the transition probability $p_{ij}(t)$ is strictly positive for \emph{all} $t > 0$, not just multiples of some period.
\end{remark}

\subsection{Recurrence and Transience of CTMCs}

\begin{definition}[Recurrence and Transience for CTMCs]
Let $(X_t) = \text{Markov}(\lambda_i, Q)$ and $\{Y_n\}$ the embedded DTMC ($Y_n = X_{J_n}$). The \textbf{first return time} to state $i$ is:
\[
T_i = \inf\{t \ge J_1 : X_t = i\}.
\]
State $i$ is \textbf{recurrent} if $\PP_i(\{t: X_t = i\} \text{ is unbounded}) = 1$, equivalently $\PP_i(T_i < \infty) = 1$ (or $q_i = 0$). State $i$ is \textbf{transient} if $q_i > 0$ and $\PP_i(T_i < \infty) < 1$.
\end{definition}

\begin{theorem}[Recurrence via Embedded DTMC]
\begin{enumerate}[label=(\roman*)]
  \item If $i$ is recurrent in $\{Y_n\}$, then $i$ is recurrent for $(X_t)$.
  \item If $i$ is transient for $\{Y_n\}$, then $i$ is transient for $(X_t)$.
  \item Every state is either recurrent or transient.
  \item Recurrence and transience are class properties.
\end{enumerate}
\end{theorem}

\begin{proof}
(i) Since $i$ is recurrent in $\{Y_n\}$, $N = \sup\{n : Y_n = i\} = \infty$ w.p.\ 1, so $X_{J_n} = i$ for infinitely many $n$. Hence $\{t : X_t = i\} \supset \{J_n : Y_n = i\}$ is unbounded w.p.\ 1.

(ii) If $i$ is transient in $\{Y_n\}$, then $N = \sup\{n : Y_n = i\} < \infty$ w.p.\ 1, so $\{t : X_t = i\} \subset [0, J_{N+1}]$, which is bounded w.p.\ 1.

(iii)–(iv) follow from the DTMC results and parts (i)–(ii).
\end{proof}

\subsection{Positive Recurrence and Stationary Distributions for CTMCs}

\begin{definition}[Positive Recurrence for CTMCs]
A recurrent state $i$ is \textbf{positive recurrent} if $m_i = \EE_i[T_i] < \infty$; otherwise it is \textbf{null recurrent}.
\end{definition}

\begin{theorem}[Stationary Distribution via Embedded DTMC]
The following are equivalent for an irreducible CTMC:
\begin{enumerate}[label=(\roman*)]
  \item $\pi$ is a stationary distribution for $Q$ (i.e., $\pi Q = 0$).
  \item $\mu P = \mu$ where $\mu_i = \pi_i q_i$ (i.e., $\mu$ is stationary for the embedded DTMC).
\end{enumerate}
\end{theorem}

\begin{proof}
$(\mu(P-I))_i = \sum_j \mu_j p_{ji} - \mu_i = \sum_j \pi_j q_j \frac{q_{ji}}{q_j} - \pi_i q_i = \sum_j \pi_j q_{ji} - \pi_i q_i = (\pi Q)_i$.
So $\mu P = \mu \Leftrightarrow \pi Q = 0$.
\end{proof}

\begin{theorem}[Positive Recurrence and Stationary Distribution for Irreducible CTMCs]
\label{thm:ctmc-posrec}
Suppose $Q$ is irreducible. The following are equivalent:
\begin{enumerate}[label=(\roman*)]
  \item Every state is positive recurrent.
  \item Some state is positive recurrent.
  \item $Q$ is non-explosive and there exists a stationary distribution $\pi$ for $Q$.
\end{enumerate}
Moreover, if (iii) holds, then $\pi_i = \dfrac{1}{m_i q_i}$.
\end{theorem}

\begin{proof}[Proof sketch]
Irreducibility implies $q_i > 0$ for all $i$.
$(i) \Rightarrow (ii)$: trivial. $(ii) \Rightarrow (iii)$: If state $i$ is positive recurrent in $\{Y_n\}$, then $(X_t)$ doesn't explode. Define:
\[
\mu_j^i = \EE_i\!\left[\int_0^{T_i} \mathbf{1}_{\{X_s = j\}}\,ds\right] = \gamma_j^i \cdot \frac{1}{q_j},
\]
where $\gamma_j^i = \EE_i\!\left[\sum_{n=0}^\infty \mathbf{1}_{\{Y_n = j,\; n < N_i\}}\right]$ counts expected visits to $j$ in the embedded DTMC before returning to $i$, and $N_i = \inf\{n \ge 1 : Y_n = i\}$.

Since $S_{n+1} \sim \Exp(q_j)$ independently, $\mu_j^i = \gamma_j^i / q_j$.

From the DTMC result, $\gamma^i P = \gamma^i$, which gives $\mu^i Q = 0$.
Moreover:
\[
\sum_j \mu_j^i = \EE_i\!\left[\int_0^{T_i} 1\,ds\right] = \EE_i[T_i] = m_i < \infty.
\]
So $\pi_j = \mu_j^i / m_i$ is a stationary distribution.

$(iii) \Rightarrow (i)$: Fix $i$. Let $v_j = \tfrac{\pi_j q_j}{\pi_i q_i}$; then $vP = v$. By the DTMC result, $v_j \ge \gamma_j^i$ componentwise, so $m_i = \sum_j \mu_j^i = \sum_j \gamma_j^i/q_j \le \sum_j v_j/q_j = \frac{1}{\pi_i q_i} < \infty$.
\end{proof}

\subsection{Convergence to Equilibrium for CTMCs}

\begin{theorem}[Convergence to Equilibrium]
Let $Q$ be irreducible and non-explosive with stationary distribution $\pi$. Then for all $i,j \in S$:
\[
  p_{ij}(t) \to \pi_j \quad \text{as } t \to \infty.
\]
(No aperiodicity condition is needed, unlike in DTMCs.)
\end{theorem}

\begin{proof}
Fix $h > 0$ and let $Z_n = X_{nh}$. Then $\{Z_n\}$ is a DTMC with transition matrix $P(h) = e^{Qh}$. Since $(X_t)$ is irreducible and non-explosive, $\{Z_n\}$ is irreducible and aperiodic (since $p_{ii}(h) > 0$ for all $h$), and $\pi$ is a stationary distribution for $P(h)$.

By the DTMC convergence theorem: $p_{ij}(nh) \to \pi_j$ as $n \to \infty$.

For general $t$, write $t = nh + r$ with $0 \le r < h$:
\[
  |p_{ij}(t) - \pi_j| \le |p_{ij}(nh) - \pi_j| + |p_{ij}(t) - p_{ij}(nh)|.
\]
The second term satisfies $|p_{ij}(t+h) - p_{ij}(t)| \le 1 - e^{-q_i h}$ (an exercise using the Kolmogorov equations). Given $\varepsilon > 0$, choose $h$ so that $1 - e^{-q_i h} < \varepsilon$, then choose $N$ large so that $|p_{ij}(nh) - \pi_j| < \varepsilon$ for $n \ge N$. Then $|p_{ij}(t) - \pi_j| < 2\varepsilon$ for all $t \ge Nh$.
\end{proof}

\subsection{Ergodic Theorem for CTMCs}

\begin{theorem}[Ergodic Theorem for CTMCs]
Let $(X_t) = \text{Markov}(\lambda, Q)$ where $Q$ is irreducible and non-explosive. Then for every bounded $f: S \to \RR$:
\[
  \frac{1}{t}\int_0^t f(X_s)\,ds \xrightarrow{a.s.} \sum_{j \in S} \pi(j)\,f(j) \quad \text{as } t \to \infty,
\]
where the right side is the ensemble (spatial) average. The result follows from the Strong Markov Property and the Strong Law of Large Numbers.
\end{theorem}

\begin{remark}
In particular, with $f = \mathbf{1}_{\{j\}}$:
\[
  \frac{1}{t}\int_0^t \mathbf{1}_{\{X_s = j\}}\,ds \xrightarrow{a.s.} \pi_j = \frac{1}{m_j q_j},
\]
which says the long-run fraction of time spent in state $j$ equals $\pi_j$.
\end{remark}

\subsection{Time Reversibility of CTMCs}

\begin{definition}[Time-Reversed CTMC]
Let $(X_t) = \text{Markov}(\pi, Q)$ where $Q$ is irreducible, non-explosive, with stationary distribution $\pi$. Fix $T > 0$ and define the \textbf{time-reversed process} $\hat{X}_t = X_{T-t}$. The reversed process has generator $\hat{Q} = (\hat{q}_{ij})$ given by:
\[
  \hat{q}_{ji} = \frac{\pi_i}{\pi_j}\,q_{ij}, \quad i \ne j.
\]
\end{definition}

\begin{theorem}[Reversed CTMC]
Let $(X_t) = \text{Markov}(\pi, Q)$ with $\pi Q = 0$. Then:
\begin{enumerate}[label=(\roman*)]
  \item $\hat{Q}$ is a valid rate matrix, irreducible and non-explosive.
  \item $\pi \hat{Q} = 0$ (i.e., $\pi$ is stationary for $\hat{Q}$).
  \item $(\hat{X}_t) = \text{Markov}(\pi, \hat{Q})$.
\end{enumerate}
\end{theorem}

\begin{proof}[Proof sketch]
For (i): $\sum_{i \ne j} \hat{q}_{ji} = \sum_{i \ne j} \frac{\pi_i}{\pi_j} q_{ij}$. Since $\pi Q = 0$, $\sum_{i} \pi_i q_{ij} = 0$, giving $\pi_j q_{jj} = -\sum_{i \ne j} \pi_i q_{ij}$, so $\sum_{i \ne j} \hat{q}_{ji} = \frac{-\pi_j q_{jj}}{\pi_j} = -q_{jj} = q_j$. Hence $\hat{q}_{jj} = -q_j$ and $\hat{Q}$ is a valid rate matrix.

For (ii): $(\pi \hat{Q})_i = \sum_j \pi_j \hat{q}_{ji} = \sum_{j \ne i} \pi_j \frac{\pi_i}{\pi_j} q_{ji} + \pi_i \hat{q}_{ii} = \pi_i \sum_{j \ne i} q_{ji} - \pi_i q_i = \pi_i(\sum_j q_{ji} - q_i) = 0$ since $\sum_j q_{ji} = 0$.

For (iii): One verifies the finite-dimensional distributions of $\hat{X}$ match those of $\text{Markov}(\pi, \hat{Q})$ by direct calculation using the stationarity of $\pi$.
\end{proof}

\begin{definition}[Reversibility for CTMCs]
$(X_t)$ started in stationarity is \textbf{reversible} w.r.t.\ $\pi$ if $(\hat{X}_t) \stackrel{d}{=} \text{Markov}(\pi, Q)$, i.e., $\hat{Q} = Q$. This is equivalent to the \textbf{detailed balance equations}:
\[
  \pi_i\,q_{ij} = \pi_j\,q_{ji}, \quad \forall\, i \ne j.
\]
\end{definition}

\begin{remark}
Detailed balance $\Rightarrow$ global balance ($\pi Q = 0$). Indeed:
\[
  (\pi Q)_i = \sum_j \pi_j q_{ji} = \sum_{j \ne i} \pi_j q_{ji} + \pi_i q_{ii} = \sum_{j \ne i} \pi_i q_{ij} - \pi_i q_i = 0.
\]
The M/M/1 queue satisfies detailed balance: $\pi_n \lambda = \pi_{n+1} \mu$ gives $\pi_{n+1} = \rho\pi_n$, yielding the geometric stationary distribution and confirming reversibility.
\end{remark}

\begin{example}[M/M/1 is Reversible]
The M/M/1 queue with $\rho < 1$ satisfies detailed balance: $\pi_n \lambda = \pi_{n+1}\mu$ for all $n \ge 0$. Hence it is time-reversible under $\pi_n = (1-\rho)\rho^n$. The reversed process is again an M/M/1 queue with the same rates.
\end{example}

\subsection{Burke's Theorem}

\begin{theorem}[Burke's Theorem]
Consider an M/M/1 queue with arrival rate $\lambda$ and service rate $\mu$, with $\rho = \lambda/\mu < 1$, started in stationarity (i.e., $X_0 \sim \pi$). Then:
\begin{enumerate}[label=(\roman*)]
  \item The departure process $D(t)$ is a Poisson process with rate $\lambda$.
  \item For each $t \ge 0$, the number of customers in the system $X_t$ is independent of the departure process $\{D(s) : 0 \le s \le t\}$.
\end{enumerate}
\end{theorem}

\begin{proof}[Proof Sketch]
Since the M/M/1 queue is time-reversible under $\pi$, the reversed process $\hat{X}_t = X_{T-t}$ is also an M/M/1 Markov chain with the same transition rates. In the reversed process, arrivals and departures interchange roles: departures in the forward process correspond to arrivals in the reversed process. Since arrivals in the reversed process form a PP($\lambda$) (by the M/M/1 input), departures in the forward process also form a PP($\lambda$). Independence of $X_t$ and past departures follows from the Markov property applied to the reversed chain.
\end{proof}

\begin{remark}
Burke's theorem is a consequence of time-reversibility. It is the key tool for analyzing networks of queues (Jackson networks), since it implies that the output of each M/M/1 node in stationarity is Poisson and independent of the node's state.
\end{remark}

\subsection{Tandem Queues}

\begin{example}[Tandem Queue]
Consider two M/M/1 queues in series: customers arrive at queue 1 at rate $\lambda$ (Poisson), are served at rate $\mu_1$, depart and immediately enter queue 2 where they are served at rate $\mu_2$. Let $\rho_k = \lambda/\mu_k < 1$ for $k=1,2$.

By Burke's theorem, the departure process from queue 1 is PP($\lambda$), independent of $X_1(t)$ (the number in queue 1). Hence the input to queue 2 is also PP($\lambda$), independent of the state of queue 1. The joint stationary distribution is:
\[
  \pi(n_1, n_2) = \pi_1(n_1)\,\pi_2(n_2)
    = (1-\rho_1)\rho_1^{n_1} \cdot (1-\rho_2)\rho_2^{n_2}, \quad n_1, n_2 \ge 0.
\]
This product-form result says the two queues in equilibrium behave as if they were independent M/M/1 queues.
\end{example}

\begin{remark}
The independence in the tandem queue is non-trivial: even though departures from queue 1 feed queue 2, the two queue lengths are independent in stationarity. This is a direct consequence of Burke's theorem and fails without time-reversibility.
\end{remark}

\subsection{Jackson Networks}

\begin{definition}[Open Jackson Network]
An \textbf{open Jackson network} consists of $K$ nodes (queues) where:
\begin{itemize}
  \item External customers arrive to node $i$ according to PP($d_i$), $d_i \ge 0$.
  \item Node $i$ has a single server with exponential service time at rate $\mu_i$.
  \item Upon service completion at node $i$, a customer routes to node $j$ with probability $p_{ij}$, or leaves the network with probability $1 - \sum_j p_{ij}$.
\end{itemize}
\end{definition}

\begin{definition}[Traffic Equations]
The \textbf{effective arrival rate} $\beta_i$ at node $i$ satisfies the \textbf{traffic equations}:
\[
  \beta_i = d_i + \sum_{j=1}^{K} \beta_j\, p_{ji}, \quad i = 1, \ldots, K.
\]
Let $\rho_i = \beta_i / \mu_i$. The network is stable if $\rho_i < 1$ for all $i$.
\end{definition}

\begin{theorem}[Existence and Uniqueness of Traffic Equation Solution]
\label{thm:jackson-traffic}
Let $P = (p_{ij})$ be the routing matrix with spectral radius $r(P) < 1$ (i.e., the network is open: every customer eventually leaves). Then the traffic equations $\beta_i = d_i + \sum_j \beta_j p_{ji}$, $i = 1,\ldots,K$, have a unique non-negative solution given by:
\[
  \boldsymbol{\beta} = \mathbf{d}(I - P)^{-1} = \mathbf{d} + \mathbf{d}P + \mathbf{d}P^2 + \cdots
\]
\end{theorem}

\begin{proof}
In matrix form the equations are $\boldsymbol{\beta} = \mathbf{d} + \boldsymbol{\beta} P$, i.e., $\boldsymbol{\beta}(I-P) = \mathbf{d}$. The matrix $(I-P)$ is invertible if and only if $r(P) < 1$ (spectral radius strictly less than 1). Since the network is open, every customer eventually leaves, so $P^n \to 0$ entry-wise as $n \to \infty$, and the Neumann series converges:
\[
  (I-P)^{-1} = \sum_{n=0}^\infty P^n.
\]
Therefore $\boldsymbol{\beta} = \mathbf{d}(I-P)^{-1}$ exists and is unique. The component $\beta_i = \sum_{n \ge 0} (\mathbf{d}P^n)_i$ counts the total rate of flow into node $i$: $(\mathbf{d}P^n)_i$ is the rate of customers who originated externally and have been routed through exactly $n$ intermediate nodes before reaching $i$. Non-negativity follows from $\mathbf{d} \ge 0$ and $P \ge 0$.
\end{proof}

\begin{remark}
The stability condition $\rho_i = \beta_i/\mu_i < 1$ for all $i$ ensures each node's queue is stable. When $\mathbf{d} = 0$ (no external arrivals), the only solution is $\boldsymbol{\beta} = 0$, confirming that the open network requires external input.
\end{remark}

\begin{theorem}[Jackson's Theorem]
For an open Jackson network in steady state with $\rho_i < 1$ for all $i$, the stationary distribution has the \textbf{product form}:
\[
  \pi(n_1, n_2, \ldots, n_K) = \prod_{i=1}^{K} \pi_i(n_i),
\]
where $\pi_i(n_i) = (1 - \rho_i)\,\rho_i^{n_i}$ is the marginal distribution of node $i$, identical to that of an M/M/1 queue with load $\rho_i$.
\end{theorem}

\begin{proof}[Proof Sketch]
The joint state space is $\mathbb{N}^K$. The joint process $(X_1(t), \ldots, X_K(t))$ is a CTMC. One verifies that $\pi(\mathbf{n}) = \prod_i \pi_i(n_i)$ satisfies the global balance equations for the network CTMC. This is done by checking that the product-form distribution satisfies the \emph{partial balance} (or local balance) equations at each node. The proof uses the fact that Poisson arrivals and exponential services give a tractable structure.
\end{proof}

\begin{remark}
Jackson's theorem is remarkable: even though queues share customers (via routing), in stationarity they are independent M/M/1 queues with effective loads $\rho_i$. This product-form property extends to networks of general $M/M/c$ nodes.
\end{remark}

\begin{example}[Jackson Network Computation]
A network with $K=2$ nodes: external arrivals $d_1 = 3$, $d_2 = 1$; routing $p_{12} = 0.5$, $p_{21} = 0.2$, all other $p_{ij} = 0$; service rates $\mu_1 = 6$, $\mu_2 = 4$.

Traffic equations:
\[
  \beta_1 = 3 + 0.2\,\beta_2, \quad \beta_2 = 1 + 0.5\,\beta_1.
\]
Solving: $\beta_2 = 1 + 0.5(3 + 0.2\beta_2) = 2.5 + 0.1\beta_2$, so $\beta_2 = 2.5/0.9 \approx 2.78$, and $\beta_1 = 3 + 0.2 \times 2.78 \approx 3.56$.

Loads: $\rho_1 = 3.56/6 \approx 0.59$, $\rho_2 = 2.78/4 \approx 0.69 < 1$. The network is stable.
\end{example}


%% ============================================================
\section{Martingales}
%% ============================================================

\subsection{Conditional Expectation}

\begin{definition}[Conditional Expectation]
Let $(\Omega, \mathcal{F}, P)$ be a probability space and $\mathcal{G} \subseteq \mathcal{F}$ a sub-$\sigma$-algebra. For an integrable random variable $X$ (i.e., $E[|X|] < \infty$), the \emph{conditional expectation} $E[X \mid \mathcal{G}]$ is the unique (a.s.) $\mathcal{G}$-measurable random variable $Z$ satisfying
\[
  \int_A Z \, dP = \int_A X \, dP \quad \text{for all } A \in \mathcal{G}.
\]
Existence and uniqueness follow from the Radon--Nikod\'ym theorem.
\end{definition}

\begin{theorem}[Properties of Conditional Expectation]
Let $X, Y$ be integrable random variables and $\mathcal{G} \subseteq \mathcal{H} \subseteq \mathcal{F}$.
\begin{enumerate}[label=(\roman*)]
  \item \textbf{Linearity:} $E[\alpha X + \beta Y \mid \mathcal{G}] = \alpha E[X \mid \mathcal{G}] + \beta E[Y \mid \mathcal{G}]$ a.s.
  \item \textbf{Tower Property:} $E\bigl[E[X \mid \mathcal{H}] \mid \mathcal{G}\bigr] = E[X \mid \mathcal{G}]$ a.s.
  \item \textbf{Measurability:} If $X$ is $\mathcal{G}$-measurable, then $E[X \mid \mathcal{G}] = X$ a.s.
  \item \textbf{Independence:} If $X$ is independent of $\mathcal{G}$, then $E[X \mid \mathcal{G}] = E[X]$ a.s.
  \item \textbf{Jensen's Inequality:} If $\phi$ is convex and $E[|\phi(X)|] < \infty$, then $\phi(E[X \mid \mathcal{G}]) \leq E[\phi(X) \mid \mathcal{G}]$ a.s.
  \item \textbf{Monotonicity:} If $X \leq Y$ a.s., then $E[X \mid \mathcal{G}] \leq E[Y \mid \mathcal{G}]$ a.s.
\end{enumerate}
\end{theorem}

\begin{remark}
The tower property is also called the \emph{law of total expectation}. When $\mathcal{G} = \{\emptyset, \Omega\}$ (trivial $\sigma$-algebra), $E[X \mid \mathcal{G}] = E[X]$. The tower property says: if you have less information ($\mathcal{G} \subseteq \mathcal{H}$), averaging again over less information gives the less-informed expectation.
\end{remark}

\begin{example}[Discrete Conditional Expectation]
Let $X$ and $Y$ be discrete random variables. Then
\[
  E[X \mid Y = y] = \sum_x x \, P(X = x \mid Y = y).
\]
In this case $E[X \mid Y]$ is the random variable $g(Y)$ where $g(y) = E[X \mid Y = y]$.

Tower property check: $E[E[X \mid Y]] = \sum_y E[X \mid Y=y] \cdot P(Y=y) = E[X]$.
\end{example}

%% ============================================================
\subsection{Filtrations and Adapted Processes}
%% ============================================================

\begin{definition}[Filtration]
A \emph{filtration} on $(\Omega, \mathcal{F}, P)$ is an increasing sequence of sub-$\sigma$-algebras
\[
  \mathcal{F}_0 \subseteq \mathcal{F}_1 \subseteq \mathcal{F}_2 \subseteq \cdots \subseteq \mathcal{F}.
\]
Intuitively, $\mathcal{F}_n$ represents the information available up to time $n$.
\end{definition}

\begin{definition}[Adapted Process]
A stochastic process $\{X_n\}_{n \geq 0}$ is \emph{adapted} to filtration $\{\mathcal{F}_n\}$ if $X_n$ is $\mathcal{F}_n$-measurable for all $n \geq 0$. That is, the value of $X_n$ is determined by the information in $\mathcal{F}_n$.
\end{definition}

\begin{definition}[Natural Filtration]
The \emph{natural filtration} generated by process $\{X_n\}$ is
\[
  \mathcal{F}_n^X = \sigma(X_0, X_1, \ldots, X_n),
\]
the smallest $\sigma$-algebra making $X_0, \ldots, X_n$ all measurable. By definition, $\{X_n\}$ is always adapted to its natural filtration.
\end{definition}

%% ============================================================
\subsection{Martingales: Definition and Examples}
%% ============================================================

\begin{definition}[Martingale]
A stochastic process $\{M_n\}_{n \geq 0}$ adapted to $\{\mathcal{F}_n\}$ with $E[|M_n|] < \infty$ for all $n$ is called:
\begin{itemize}
  \item A \emph{martingale} if $E[M_{n+1} \mid \mathcal{F}_n] = M_n$ a.s. for all $n$.
  \item A \emph{submartingale} if $E[M_{n+1} \mid \mathcal{F}_n] \geq M_n$ a.s. for all $n$.
  \item A \emph{supermartingale} if $E[M_{n+1} \mid \mathcal{F}_n] \leq M_n$ a.s. for all $n$.
\end{itemize}
\end{definition}

\begin{remark}
Intuitively: a martingale models a \emph{fair game} (conditional expected future value equals present value). A submartingale models a \emph{favorable} game (you expect to gain on average). A supermartingale models an \emph{unfavorable} game.

Note: if $\{M_n\}$ is a martingale, then $E[M_n] = E[M_0]$ for all $n$ (by the tower property applied repeatedly).
\end{remark}

\begin{example}[Simple Random Walk]
Let $X_1, X_2, \ldots$ be i.i.d.\ with $P(X_i = +1) = P(X_i = -1) = 1/2$. Define $S_n = X_1 + \cdots + X_n$ and $\mathcal{F}_n = \sigma(X_1, \ldots, X_n)$. Then:
\[
  E[S_{n+1} \mid \mathcal{F}_n] = E[S_n + X_{n+1} \mid \mathcal{F}_n] = S_n + E[X_{n+1}] = S_n + 0 = S_n.
\]
So $\{S_n\}$ is a martingale with respect to $\{\mathcal{F}_n\}$.
\end{example}

\begin{example}[Product Martingale]
Let $X_1, X_2, \ldots$ be i.i.d.\ with $E[X_i] = 1$ (e.g., $P(X_i = 2) = P(X_i = 0) = 1/2$). Define $M_n = X_1 X_2 \cdots X_n$, $M_0 = 1$. Then:
\[
  E[M_{n+1} \mid \mathcal{F}_n] = M_n \cdot E[X_{n+1}] = M_n \cdot 1 = M_n.
\]
So $\{M_n\}$ is a martingale.
\end{example}

\begin{example}[Doob Martingale]
Let $Z$ be an integrable random variable and $\{\mathcal{F}_n\}$ any filtration. Define
\[
  M_n = E[Z \mid \mathcal{F}_n].
\]
Then $\{M_n\}$ is a martingale: by the tower property,
\[
  E[M_{n+1} \mid \mathcal{F}_n] = E[E[Z \mid \mathcal{F}_{n+1}] \mid \mathcal{F}_n] = E[Z \mid \mathcal{F}_n] = M_n.
\]
This is called the \emph{Doob martingale} or \emph{Lévy's upward theorem} construction.
\end{example}

\begin{example}[Exponential Martingale]
For a simple random walk $S_n = \sum_{i=1}^n X_i$ with $X_i \in \{-1, +1\}$ equally likely, define for $\theta \in \mathbb{R}$:
\[
  M_n = \frac{e^{\theta S_n}}{(E[e^{\theta X_1}])^n} = \frac{e^{\theta S_n}}{\left(\cosh\theta\right)^n}.
\]
One can verify $E[M_{n+1} \mid \mathcal{F}_n] = M_n$, so this is a martingale. Exponential martingales are key tools for obtaining concentration inequalities.
\end{example}

%% ============================================================
\subsection{Stopping Times}
%% ============================================================

\begin{definition}[Stopping Time]
A random variable $\tau : \Omega \to \{0, 1, 2, \ldots\} \cup \{\infty\}$ is a \emph{stopping time} with respect to filtration $\{\mathcal{F}_n\}$ if
\[
  \{\tau \leq n\} \in \mathcal{F}_n \quad \text{for all } n \geq 0.
\]
Equivalently, $\{\tau = n\} \in \mathcal{F}_n$ for all $n$. The event ``we stop at or before time $n$'' must be determined by information available at time $n$ alone.
\end{definition}

\begin{remark}
Examples of stopping times: the first time a process hits a set $A$ (hitting time $\tau_A = \inf\{n : X_n \in A\}$), a fixed deterministic time $\tau = N$, and the minimum of two stopping times.

A \emph{non-example}: ``the last time $S_n = 0$ before time $N$'' is NOT a stopping time, because determining it requires future information.
\end{remark}

\begin{definition}[Stopped Process and $\sigma$-algebra $\mathcal{F}_\tau$]
For a stopping time $\tau$, the \emph{stopped process} is $M_{\tau \wedge n} = M_{\min(\tau, n)}$.

The \emph{$\sigma$-algebra at stopping time $\tau$} is
\[
  \mathcal{F}_\tau = \{A \in \mathcal{F} : A \cap \{\tau \leq n\} \in \mathcal{F}_n \text{ for all } n\}.
\]
$\mathcal{F}_\tau$ represents the information accumulated up to the (random) stopping time.
\end{definition}

\begin{lemma}[Stopped Martingale]
If $\{M_n\}$ is a martingale and $\tau$ is a stopping time, then the stopped process $\{M_{\tau \wedge n}\}$ is also a martingale.
\end{lemma}

\begin{proof}
By adaptedness and the martingale property:
\[
  E[M_{\tau \wedge (n+1)} \mid \mathcal{F}_n]
  = M_{\tau \wedge n} \cdot \mathbf{1}_{\tau \leq n} + E[M_{n+1} \mid \mathcal{F}_n] \cdot \mathbf{1}_{\tau > n}
  = M_{\tau \wedge n} \cdot \mathbf{1}_{\tau \leq n} + M_n \cdot \mathbf{1}_{\tau > n}
  = M_{\tau \wedge n}.
\]
\end{proof}

%% ============================================================
\subsection{Optional Stopping Theorem}
%% ============================================================

\begin{theorem}[Optional Stopping Theorem (Doob's OST)]
Let $\{M_n\}$ be a martingale and $\tau$ a stopping time. Then $E[M_\tau] = E[M_0]$ provided any one of the following conditions holds:
\begin{enumerate}[label=(\roman*)]
  \item $\tau$ is bounded: $\tau \leq N$ a.s.\ for some fixed $N < \infty$;
  \item $\{M_n\}$ is uniformly integrable and $\tau < \infty$ a.s.;
  \item $E[\tau] < \infty$ and $|M_{n+1} - M_n| \leq c$ a.s.\ (bounded increments).
\end{enumerate}
\end{theorem}

\begin{proof}[Proof (bounded case)]
Since $\tau \leq N$, we have $M_\tau = M_{\tau \wedge N}$. Since the stopped process is a martingale:
\[
  E[M_\tau] = E[M_{\tau \wedge N}] = E[M_0].
\]
The last equality uses the martingale property of $M_{\tau \wedge n}$ iterated from $n=0$ to $n=N$.
\end{proof}

\begin{remark}
The bounded case follows simply because we can write:
\[
  M_\tau = M_0 + \sum_{k=1}^{N} \mathbf{1}_{\tau \geq k}(M_k - M_{k-1}).
\]
Taking expectations, each term $E[\mathbf{1}_{\tau \geq k}(M_k - M_{k-1})] = 0$ because $\{\tau \geq k\} = \{\tau \leq k-1\}^c \in \mathcal{F}_{k-1}$ and $E[M_k - M_{k-1} \mid \mathcal{F}_{k-1}] = 0$.
\end{remark}

\begin{example}[Gambler's Ruin via OST]
Let $S_n = X_1 + \cdots + X_n$ be a simple random walk with $P(X_i = +1) = P(X_i = -1) = 1/2$. The gambler starts at position $s \in (0, a+b)$ and plays until hitting $0$ (ruin) or $a+b$ (win), where the barrier positions are $0$ and $a+b$.

More precisely: start at $S_0 = b$ (positive integer), stop at $\tau = \inf\{n : S_n = 0 \text{ or } S_n = a+b\}$. Let $p$ = probability of reaching $a+b$ before $0$.

\textbf{Method 1:} $\{S_n\}$ is a martingale. Applying OST (condition (i) or (iii)):
\[
  E[S_\tau] = E[S_0] = b.
\]
But $S_\tau \in \{0, a+b\}$, so:
\[
  E[S_\tau] = (a+b) \cdot p + 0 \cdot (1-p) = (a+b) \cdot p.
\]
Therefore $p = b/(a+b)$.

\textbf{Method 2:} $\{S_n^2 - n\}$ is a martingale (since $E[S_{n+1}^2 \mid \mathcal{F}_n] = S_n^2 + 1$). Applying OST:
\[
  E[S_\tau^2 - \tau] = E[S_0^2 - 0] = b^2.
\]
Since $E[S_\tau^2] = (a+b)^2 \cdot p = (a+b)^2 \cdot \frac{b}{a+b} = b(a+b)$, we get $E[\tau] = b(a+b) - b^2 = ab$.
\end{example}

%% ============================================================
\subsection{Martingale Convergence}
%% ============================================================

\begin{theorem}[Martingale Convergence Theorem]
Let $\{M_n\}$ be a martingale (or submartingale) with $\sup_n E[|M_n|] < \infty$. Then there exists a random variable $M_\infty$ with $E[|M_\infty|] < \infty$ such that
\[
  M_n \xrightarrow{a.s.} M_\infty \quad \text{as } n \to \infty.
\]
\end{theorem}

\begin{remark}
Note that a.s.\ convergence does NOT necessarily imply $L^1$ convergence. For $L^1$ convergence (i.e., $E[|M_n - M_\infty|] \to 0$), we need the additional condition of \emph{uniform integrability}.
\end{remark}

\begin{definition}[Uniform Integrability]
A family $\{X_\alpha\}$ of random variables is \emph{uniformly integrable (UI)} if
\[
  \lim_{K \to \infty} \sup_\alpha E\bigl[|X_\alpha| \mathbf{1}_{|X_\alpha| > K}\bigr] = 0.
\]
Equivalently: the tails of all $X_\alpha$ are uniformly small in expectation.
\end{definition}

\begin{theorem}[$L^1$ Martingale Convergence]
A martingale $\{M_n\}$ converges in $L^1$ (and hence a.s.) if and only if it is uniformly integrable. In this case, $M_\infty = \lim_n M_n$ exists a.s.\ and $E[M_n \mid \mathcal{F}_m] = M_m$ for all $m \leq n$ extends to $E[M_\infty \mid \mathcal{F}_n] = M_n$.
\end{theorem}

\begin{example}[Galton--Watson Branching Process and Martingale]
\label{ex:galton-watson}
\textbf{Setup.} A population evolves in generations. Each individual in generation $n$ independently produces a random number of offspring according to a distribution with pgf $G(s) = \EE[s^L]$ where $L$ counts offspring. Let $Z_0 = 1$, $Z_n = $ size of generation $n$, and $\mu = \EE[L] = G'(1)$.

\textbf{The martingale.} Define $M_n = Z_n / \mu^n$ (normalized population size). Then $\{M_n\}$ is a non-negative martingale with respect to $\mathcal{F}_n = \sigma(Z_0, Z_1, \ldots, Z_n)$:
\[
  \EE[M_{n+1} \mid \mathcal{F}_n] = \frac{1}{\mu^{n+1}} \EE[Z_{n+1} \mid Z_n] = \frac{1}{\mu^{n+1}} \cdot \mu Z_n = \frac{Z_n}{\mu^n} = M_n.
\]
Here we used $\EE[Z_{n+1} \mid Z_n = k] = k\mu$ (each of $k$ individuals produces mean $\mu$ offspring independently).

\textbf{Convergence.} Since $M_n \ge 0$ and $\EE[M_n] = \EE[M_0] = 1 < \infty$, by the Martingale Convergence Theorem, $M_n \to M_\infty$ a.s.

\textbf{Extinction probability.} Let $q = \PP(Z_n = 0 \text{ for some } n) = \PP(\text{extinction})$. Then $q$ is the smallest non-negative solution to $q = G(q)$:
\begin{itemize}
  \item If $\mu = \EE[L] \le 1$ (subcritical or critical): $q = 1$ (certain extinction).
  \item If $\mu > 1$ (supercritical): $q < 1$ (survival with positive probability $1-q$).
\end{itemize}
On the extinction event, $M_\infty = 0$. On survival ($\mu > 1$), $M_\infty > 0$ with positive probability, capturing the limiting normalized population size.

\textbf{Second moment and $L^2$ convergence.} If $\sigma^2 = \text{Var}(L) < \infty$, then $\EE[M_n^2]$ stays bounded (compute via $\EE[Z_n^2]$ recursion), so $M_n \to M_\infty$ in $L^2$ as well. In particular, $\EE[M_\infty] = 1 \cdot \PP(\text{survival}) + 0 \cdot \PP(\text{extinction})$, which with $L^2$ convergence gives $\EE[M_\infty] = 1-q$ only when $\mu > 1$ and $\sigma^2 < \infty$.
\end{example}

%% ============================================================
\subsection{Azuma--Hoeffding Inequality}
%% ============================================================

\begin{definition}[Martingale Difference Sequence]
Given a martingale $\{M_n\}$ with respect to $\{\mathcal{F}_n\}$, the \emph{martingale differences} are $D_k = M_k - M_{k-1}$ for $k \geq 1$. Note that $E[D_k \mid \mathcal{F}_{k-1}] = 0$ and $M_n = M_0 + \sum_{k=1}^n D_k$.
\end{definition}

\begin{definition}[Bounded Differences Condition]
A martingale $\{M_n\}$ satisfies the \emph{bounded differences condition} with constants $c_1, c_2, \ldots > 0$ if
\[
  |D_k| = |M_k - M_{k-1}| \leq c_k \quad \text{a.s.\ for all } k \geq 1.
\]
\end{definition}

\begin{theorem}[Azuma--Hoeffding Inequality]
Let $\{M_n\}_{n=0}^N$ be a martingale satisfying the bounded differences condition with constants $c_1, \ldots, c_N$. Then for any $t > 0$:
\[
  P(M_N - M_0 \geq t) \leq \exp\!\left(-\frac{t^2}{2 \sum_{k=1}^N c_k^2}\right),
\]
and similarly:
\[
  P(|M_N - M_0| \geq t) \leq 2\exp\!\left(-\frac{t^2}{2 \sum_{k=1}^N c_k^2}\right).
\]
\end{theorem}

\begin{proof}[Proof Sketch]
The proof uses the Hoeffding lemma: if $E[D] = 0$ and $|D| \leq c$ a.s., then $E[e^{\lambda D}] \leq e^{\lambda^2 c^2 / 2}$.

Define $L_n = e^{\lambda M_n} / \prod_{k=1}^n E[e^{\lambda D_k} \mid \mathcal{F}_{k-1}]$. Using the Hoeffding lemma and Markov's inequality on the exponential martingale $\{e^{\lambda M_n}\}$, we get:
\[
  P(M_N - M_0 \geq t) = P\!\left(e^{\lambda(M_N - M_0)} \geq e^{\lambda t}\right) \leq e^{-\lambda t} \cdot E[e^{\lambda(M_N - M_0)}] \leq e^{-\lambda t + \lambda^2 \sum c_k^2 / 2}.
\]
Optimizing over $\lambda > 0$ (set $\lambda = t / \sum c_k^2$) gives the result.
\end{proof}

\begin{corollary}[McDiarmid's Inequality]
Let $f(x_1, \ldots, x_n)$ be a function of $n$ independent random variables satisfying the bounded differences condition:
\[
  \sup_{x_1,\ldots,x_n,\,x_i'} |f(\ldots, x_i, \ldots) - f(\ldots, x_i', \ldots)| \leq c_i.
\]
Then with $Z = f(X_1, \ldots, X_n)$ and $\mu = E[Z]$:
\[
  P(|Z - \mu| \geq t) \leq 2 \exp\!\left(-\frac{2t^2}{\sum_{i=1}^n c_i^2}\right).
\]
This follows from Azuma--Hoeffding applied to the Doob martingale $M_k = E[Z \mid X_1, \ldots, X_k]$.
\end{corollary}

\begin{example}[Azuma--Hoeffding Application]
Consider a simple random walk $S_n = X_1 + \cdots + X_n$ with $X_i \in \{-1, +1\}$. The differences are $|S_n - S_{n-1}| = 1$, so $c_k = 1$ for all $k$. By Azuma--Hoeffding:
\[
  P(S_n \geq t) \leq e^{-t^2/(2n)}.
\]
Equivalently, $P(S_n \geq t\sqrt{n}) \leq e^{-t^2/2}$, consistent with the CLT approximation $S_n / \sqrt{n} \approx N(0,1)$.
\end{example}

\begin{example}[Hamming Distance via Azuma--Hoeffding (PS13 type)]
Let $X_1, \ldots, X_n \in \{0,1\}$ be i.i.d.\ Bernoulli$(p)$ bits (the transmitted codeword), and let $f(X_1,\ldots,X_n)$ be the Hamming distance from the received word to the nearest codeword in some fixed code $\mathcal{C}$:
\[
  f(x_1,\ldots,x_n) = \min_{c \in \mathcal{C}} d_H(x, c), \quad d_H(x,c) = \sum_{i=1}^n \mathbf{1}[x_i \neq c_i].
\]
Changing any single bit $x_i \to x_i'$ changes the Hamming distance by at most 1, so $f$ satisfies the bounded differences condition with $c_i = 1$ for all $i$. By McDiarmid's inequality:
\[
  P\bigl(|f(X_1,\ldots,X_n) - \EE[f]| \geq t\bigr) \leq 2e^{-2t^2/n}.
\]
This says the minimum distance decoding error is concentrated around its mean: with high probability, $f$ is within $O(\sqrt{n \log(1/\delta)})$ of its expectation.
\end{example}

\begin{example}[Graph Chromatic Number / Shortest Path via Azuma--Hoeffding (PS13 type)]
\textbf{Chromatic number:} Let $G(n,p)$ be the Erd\H{o}s--R\'enyi random graph on $n$ vertices. Let $\chi(G)$ be the chromatic number. Reveal the edges one by one; revealing edge $\{i,j\}$ changes $\chi$ by at most 1. There are $\binom{n}{2}$ edges, each bounded by $c_k = 1$. By Azuma--Hoeffding:
\[
  P\bigl(|\chi(G) - \EE[\chi(G)]| \geq t\bigr) \leq 2\exp\!\left(-\frac{2t^2}{\binom{n}{2}}\right) = 2e^{-4t^2/n(n-1)}.
\]

\textbf{Shortest path:} Let $d(s,t)$ be the shortest path length between two fixed vertices $s,t$ in $G(n,p)$. Changing one edge changes $d(s,t)$ by at most 1, so:
\[
  P\bigl(|d(s,t) - \EE[d(s,t)]| \geq t\bigr) \leq 2e^{-4t^2/n(n-1)}.
\]
These examples illustrate how Azuma--Hoeffding via the Doob martingale $M_k = \EE[f \mid \text{first } k \text{ edges}]$ gives concentration for any Lipschitz function of random inputs.
\end{example}

%% ============================================================
\subsection{Wald's Identity}
%% ============================================================

\begin{theorem}[Wald's Identity]
Let $X_1, X_2, \ldots$ be i.i.d.\ with $E[X_1] = \mu < \infty$, and let $\tau$ be a stopping time with respect to the natural filtration $\mathcal{F}_n = \sigma(X_1,\ldots,X_n)$, with $E[\tau] < \infty$. Define $S_n = X_1 + \cdots + X_n$. Then
\[
  E[S_\tau] = \mu \cdot E[\tau].
\]

\textbf{Conditions required for Wald's identity to hold:}
\begin{enumerate}[label=(\roman*)]
  \item $X_1, X_2, \ldots$ are \textbf{i.i.d.} (identically distributed and mutually independent).
  \item $E[|X_1|] < \infty$ (finite mean $\mu$).
  \item $\tau$ is a \textbf{stopping time} — i.e., $\{\tau = n\} \in \mathcal{F}_n = \sigma(X_1,\ldots,X_n)$ for all $n$. Crucially, $\tau$ may depend on $X_1, X_2, \ldots$ but only through their history, not future values.
  \item $E[\tau] < \infty$.
\end{enumerate}
\end{theorem}

\begin{proof}
Write $S_\tau = \sum_{k=1}^\infty X_k \mathbf{1}_{k \leq \tau}$. Then:
\[
  E[S_\tau] = \sum_{k=1}^\infty E[X_k \mathbf{1}_{k \leq \tau}].
\]
Since $\{\tau \geq k\} = \{\tau \leq k-1\}^c \in \mathcal{F}_{k-1}$ and $X_k$ is independent of $\mathcal{F}_{k-1}$:
\[
  E[X_k \mathbf{1}_{k \leq \tau}] = E[X_k] \cdot P(\tau \geq k) = \mu \cdot P(\tau \geq k).
\]
Summing: $E[S_\tau] = \mu \sum_{k=1}^\infty P(\tau \geq k) = \mu \cdot E[\tau]$.
\end{proof}

\begin{remark}
Alternatively: $M_n = S_n - n\mu$ is a martingale. By the Optional Stopping Theorem (under $E[\tau] < \infty$ and bounded increments),
\[
  E[M_\tau] = E[S_\tau - \tau\mu] = E[M_0] = 0 \implies E[S_\tau] = \mu E[\tau].
\]
\end{remark}

\begin{example}[Counterexample: Wald fails when $\tau$ is not a stopping time]
\label{ex:wald-fail}
Let $X_1, X_2, \ldots$ be i.i.d.\ with $P(X_i = +1) = P(X_i = -1) = 1/2$, so $\mu = 0$.

\textbf{Failure case 1: $\tau$ depends on future values.} Define $\tau = \arg\max_{1 \le k \le n} S_k$ (the time index at which the maximum is achieved among the first $n$ steps). This is NOT a stopping time since determining whether $k$ is the maximizer requires knowledge of $X_{k+1}, \ldots, X_n$ (future information). Wald's formula gives $E[S_\tau] = \mu \cdot E[\tau] = 0$, but in fact $E[S_\tau] = E[\max_{k \le n} S_k] > 0$ for $n \ge 1$.

\textbf{Failure case 2: $E[\tau] = \infty$.} Let $X_i$ be i.i.d.\ with $P(X_i = 2) = P(X_i = -1) = 1/2$, so $\mu = 1/2$. Let $\tau = \min\{n : S_n \le 0\}$. One can show $E[\tau] = \infty$ in this case (because positive drift makes return to 0 very slow). Then $E[S_\tau] = E[S_\tau \mid S_\tau \le 0] \le 0$, but $\mu \cdot E[\tau] = \infty$, so Wald's identity fails.

\textbf{Failure case 3 (PS13 Q2 type): $\tau$ not independent of $X$'s through stopping time structure.} Suppose we define $\tau = n$ if $X_n = 1$ and $\tau = 1$ otherwise (i.e., $\tau$ is chosen to land on a favorable observation). This is not a stopping time (the choice of $\tau$ depends on $X_\tau$, which is future information at time $\tau - 1$). Computing directly: $E[S_\tau] \ne \mu \cdot E[\tau]$ in general.

\textbf{Key takeaway:} Wald's identity requires $\tau$ to be a genuine stopping time — determined by the past — and $E[\tau] < \infty$. If either condition fails, the identity can break down.
\end{example}

\begin{example}[Expected Stopping Time for Random Walk]
For the simple random walk (with $\mu = 0$) starting at $S_0 = b$ and stopping at $\{0, a+b\}$:

From the gambler's ruin calculation, $E[\tau] = ab$ (derived using $M_n = S_n^2 - n$ and OST).

As another application: for i.i.d.\ $X_i \in \{-1, +1\}$ with $P(X_i = +1) = p \neq 1/2$, define $\mu = 2p-1 \neq 0$. Then $S_n - n\mu$ is a martingale. By Wald: if we stop at $\tau = \inf\{n : S_n \notin (0, N)\}$, then $E[S_\tau] = b + \mu \cdot E[\tau]$, giving $E[\tau]$ in terms of the boundary hitting probabilities.
\end{example}

%% ============================================================
\subsection{Applications to Queueing}
%% ============================================================

\begin{remark}[Martingales in Queueing Theory]
Martingale methods provide powerful tools for analyzing random processes in queueing:
\begin{itemize}
  \item \textbf{Stability:} If $Q_n$ is the queue length process, $\{Q_n - \rho n\}$ for appropriate $\rho$ can be a supermartingale, establishing stability conditions via convergence theorems.
  \item \textbf{Waiting time bounds:} Azuma--Hoeffding applied to service time sums yields concentration bounds on waiting times.
  \item \textbf{Fluid models:} Doob martingales arise naturally in the analysis of fluid-flow approximations to queues.
\end{itemize}
\end{remark}

\begin{example}[Martingale Analysis of M/G/1 Queue]
In an M/G/1 queue with Poisson arrivals (rate $\lambda$), service times $B_i$ (i.i.d., mean $1/\mu$, second moment $E[B^2]$), let $Q_n$ be the queue length just after the $n$-th departure.

The sequence $\{Q_n\}$ is a Markov chain. Define $\rho = \lambda/\mu < 1$ (stability condition). For the stationary analysis, a key martingale identity is the Pollaczek--Khinchine (P-K) formula:
\[
  E[Q] = \rho + \frac{\rho^2 + \lambda^2 E[B^2]}{2(1-\rho)}.
\]
The mean waiting time in system is $W = E[Q]/\lambda + 1/\mu$ (Little's Law).
\end{example}

%% ============================================================
\subsection{Summary of Key Results}
%% ============================================================

\begin{remark}[Course Summary]
The course IE622 covered the following major topics:
\begin{enumerate}[label=(\roman*)]
  \item \textbf{Renewal Theory:} Renewal processes, renewal reward theorem, key renewal theorem, alternating renewal processes.
  \item \textbf{Markov Chains (DTMC):} Chapman--Kolmogorov equations, classification of states (transient/recurrent/positive recurrent), stationary distributions, ergodic theorem, birth-death chains.
  \item \textbf{Markov Chains (CTMC):} Generator matrices, Kolmogorov forward/backward equations, stationary distributions, PASTA property, birth-death processes.
  \item \textbf{Queueing Theory:} M/M/1, M/M/c, M/M/$\infty$ queues, Little's Law, Burke's theorem, tandem queues, Jackson networks with product-form solutions.
  \item \textbf{Martingale Theory:} Conditional expectation, filtrations, martingales/submartingales/supermartingales, stopping times, Optional Stopping Theorem, martingale convergence, Azuma--Hoeffding inequality, Wald's identity.
\end{enumerate}
The unifying theme is the analysis of \emph{stochastic processes} — random systems evolving over time — using tools from probability theory to derive exact or approximate performance measures.
\end{remark}

\end{document}
